%% GENERATED FILE -- do not edit by hand.
%% Regenerate with: python scripts/build_blinded_submission.py
%% Source: real-patterns-need-closure-SYNTHESE-sn.tex + ../../real-patterns-need-closure-SYNTHESE-mainbody.tex
%% Identifying metadata is deliberately omitted here; it lives in title-page.tex.
\documentclass[pdflatex,sn-basic]{sn-jnl}

%%%% Standard Packages
\usepackage[T1]{fontenc}
\usepackage[utf8]{inputenc}
\usepackage{lmodern}
\usepackage{graphicx}%
\usepackage{amsmath,amssymb,amsfonts}%
\usepackage{amsthm}%
\usepackage{xcolor}%
\usepackage{textcomp}%
\usepackage{booktabs}%
\providecommand{\tightlist}{%
  \setlength{\itemsep}{0pt}\setlength{\parskip}{0pt}}

\raggedbottom

\begin{document}

\title[Operationalizing Screening Off]{Operationalizing Screening Off for Rainforest Admission}

\abstract{Scale-relative structural realism allows that real patterns may occupy more than one level of ontology, but the admittance criteria for those levels remain underdetermined. Franklin and Robertson (2024) offer the most developed recent proposal: emergence, understood as screening off plus novelty, earns an entity its place in the rainforest. Their criterion is the right one, but its stated test underdetermines it. Conditional independence admits an observational reading, on which micro-detail is redundant across the system's actual trajectory, and a modal reading, on which that redundancy survives admissible intervention. They coincide in the central cases and diverge where apparent macro-autonomy is borrowed from off-partition structure. A macro-variable driven by a shared upstream process can satisfy their conditions in full, novelty included, while distinct micro-realizations of its macrostate come apart under perturbation. On the modal reading, screening off and transition autonomy are not two conditions but one: to screen off micro-detail from the macro-future under admissible perturbation is to carry autonomous transition structure at that grain. Closure is the admissibility-indexed test that recovers this modal content. A candidate qualifies when, for a fixed regime, horizon, and admissible intervention class, macrostate information suffices for macro-transitions. Strong lumpability is the exact benchmark; non-ideal cases are graded by leakiness, which must itself be indexed to the intervention class rather than to the observed distribution. Interventions serve throughout as diagnostic probes rather than as metaphysical primitives. The contribution is a correction to how an existing filter is read and applied rather than a rival to it.}

\keywords{real patterns, rainforest realism, closure, strong lumpability, structural realism, interventionism}

\maketitle

\section{Introduction: Closure and the Admittance
  Problem}\label{introduction-closure-and-the-admittance-problem}

Lewis observed that we have no principled reason to deny existence to
the mereological sum of ``the right half of my left shoe plus the Moon
plus the sum of all Her Majesty's ear-rings,'' however sensible it is
to ignore such things in ordinary thought \citep[p.~213]{lewis1986}.
Fair enough. But consider the difference between that composite and a
hurricane. Knowing a hurricane's pressure organization and rotational
structure tells you what it will do next without tracking every
molecule. The shoe-moon composite carries no comparable transition
structure; predicting its future requires independently tracking each
component. Grant that both exist. What remains open is which patterns
carry autonomous macro-dynamics and therefore deserve robust
commitment.

Scale-relative structural realism says that the world is not exhausted
by a single privileged descriptive level. The special sciences succeed
by tracking higher-level patterns stable enough to support prediction,
explanation, and intervention, and Ladyman and Ross's rainforest realism
gave this thought its most influential metaphysical form: reality is
populated by real patterns at many scales, not only by microphysical
furniture \citep{ladyman2007}. But one question remains insufficiently
settled. Which patterns earn admission into that rainforest, and which
should be treated more cautiously as useful summaries, temporary
proxies, or dynamically idle aggregates?

Franklin and Robertson offer the most developed recent answer
\citep{franklinrobertson2024}. Emergence earns rainforest admission
when the macro-variable renders micro-details conditionally irrelevant
(screening off) relative to the explanandum, and the macro-variable
participates in a dependency profile that is not type-identical to the
realizing micro-dependencies. That is a substantial advance. It
captures why obvious gerrymanders and trout-turkey style composites
should be excluded, while preserving the thought that higher-level
ontology is earned by stable macro-level dependence rather than by
intuitive naturalness alone. They explicitly cast the issue as one of
the ``admittance criteria for the rainforest''
\citep[\S1]{franklinrobertson2024}, and they do so in a way meant to
remain compatible with theoretical reducibility rather than to revive a
strong irreducibility test.

Screening off plus novelty is indispensable, but admits two readings of
conditional independence, observational and modal. No published test in
this literature fixes the modal reading, and the two readings return
opposite verdicts on candidates the debate cares about. A candidate
macro-variable deserves robust commitment only under the modal reading: it must carry
its own transition structure under a fixed regime, horizon, and
admissible intervention class. In exact Markov settings, strong
lumpability is the benchmark realization of that condition. In
non-ideal settings, closure is graded via leakiness and convergent
diagnostics under fixed constraints, with explicit downgrade rules when
verdicts are fragile. A candidate can satisfy observational
screening-off conditions and enter into stable, law-like
macrodependencies while still borrowing its apparent autonomy from
upstream structure or from hidden micro-identity. Closure, indexed to
admissible interventions, is the test that decides between the two
readings.

The relation to earlier work is cumulative. Dennett's real patterns
identify candidate structures through compression and projectibility \citep{dennett1991}; Franklin and
Robertson supply the admittance filter by requiring screening off plus
novelty \citep{franklinrobertson2024}; the criterion developed here
specifies which reading of that filter earns admission and provides
the test that discriminates it. The candidates it withholds commitment
from are those that pass the filter observationally while remaining
parasitic on off-partition structure, whether hidden micro-identity, a
shared upstream cause, or a fast hidden relaxation process. The aim is
to sharpen rainforest realism from inside its own commitments rather
than to displace them.

The wedge that forces the further step is a class of parasitism cases
in which observational success masks off-partition dependence. The
canonical example is the latent-driver proxy: a macro-variable
aggregating two subsystems driven by a shared upstream process can
satisfy observational screening off because the driver is a common
cause of the proxy's present and its future, so that the proxy's
current state carries exactly the upstream information its next state
turns on. Under passive observation the proxy looks
transition-autonomous. Perturb the proxy directly with an admissible
intervention while leaving the drivers intact, and hidden
heterogeneity reappears: within-class micro-differences now matter for
the ensuing macro-transition.
Section~\ref{what-the-formal-machinery-alone-does-not-provide} makes
this case explicit as a countermodel rather than an intuition, since
the argument turns on the claim that a candidate can satisfy the
stated admittance conditions in full and still fail transition
autonomy. A
second case, hidden-feedback parasitism
(\S\ref{hidden-feedback-case}), shows the same failure signature from a
different source, when fast off-partition dynamics rather than a shared
upstream driver sustain the apparent autonomy.

Throughout this paper, \emph{transition autonomy} names the property
the criterion tests: the candidate partition's own dynamics close under
admissible perturbation, so within-class micro-detail adds no
predictive power. \emph{Gerrymandered} then means dynamically
incoherent: a partition that fails transition autonomy by requiring
persistent within-class micro-bookkeeping, not merely one that is
intuitively disjunctive or aesthetically odd. \emph{Parasitic} means
that a macro-variable's apparent autonomy is borrowed from upstream
structure rather than realized in the candidate partition's own
transition organization.

Three claims of different strength are in play, and readers who grant
one but not the next can take the paper's contribution in corresponding
part. The weakest is that closure is a necessary anti-gerrymandering
constraint on serious macro-ontology claims. The middle claim is that,
given the structural realist and interventionist commitments stated
below, closure under admissible intervention is the appropriate
criterion for robust macro-level commitment. The strongest is
programmatic: rainforest realism is naturally read as a structured
space of computationally closed coarse-grainings rather than as a
binary inventory of what is, or is not, real. Sections~\ref{closure-and-admissibility}
through~\ref{approximate-closure-without-instrumentalist-drift} develop
the first two claims; the structured-space reading is offered as a
programmatic upshot in
Section~\ref{conclusion-a-disciplined-realism-claim}, with the formal
richness of that space in any given domain to be cashed out as the
closure discipline matures.

The paper is a philosophical argument. It provides a criterion and a
disciplined protocol for applying it; estimator selection, finite-sample behavior, and domain-specific
implementation remain downstream methodological tasks.

\subsection{Novelty and Positioning}\label{novelty-and-positioning}

The novelty claim has three parts.

\begin{enumerate}
	\item
	      Recovering screening off's modal content: the contribution
	      corrects how screening off is read; it does not add a second
	      condition on top of it. Conditional
	      independence admits an observational reading, on which
	      micro-detail is redundant given the macrostate across the
	      system's actual trajectory, and a modal reading, on which that
	      redundancy survives admissible intervention. These readings
	      coincide in the central cases and come apart exactly where
	      apparent macro-autonomy is borrowed from off-partition
	      structure. Franklin and Robertson's criterion is retained in
	      full; the claim is that their stated test, like the
	      screening-off diagnostics it draws on, fixes only the
	      observational reading, whereas the reading that carries
	      ontological weight is the modal one. On the modal reading,
	      screening off and transition autonomy are the same property.
	      To screen off micro-detail from the macro-future under
	      admissible perturbation just is to carry autonomous transition
	      structure at that grain. Closure is the admissibility-indexed
	      test that recovers this modal content and thereby distinguishes
	      autonomous macro-variables from parasitic ones. Their account
	      is valuable precisely because it frames rainforest admission
	      without treating reducibility by itself as disqualifying
	      \citep{franklinrobertson2024}. The observational reading also
	      appears outside their paper. Meyer generalizes the same package beyond the
	      special sciences and states the point at issue explicitly,
	      defining dependence as unconditional relevance and
	      \emph{autonomy} as conditional irrelevance,
	      \(P(E_2 \mid E_1 \,\&\, B) \approx P(E_2 \mid E_1)\)
	      \citep[p.~19]{meyer2026}. Autonomy is there identified with a
	      probabilistic condition on an unindexed distribution, which is
	      the identification this paper argues underdetermines
	      what it is meant to secure.
	\item
	      Beyond formal closure results alone: recent work distinguishes
	      informational, causal, and computational closure as formal
	      diagnostics within the formal settings studied
	      \citep{rosas2024}, but does not say when those
	      diagnostics warrant ontological commitment or when commitment
	      should be withdrawn. Those diagnostics are also defined
	      over the observed process, and their authors are explicit that
	      the counterfactual reading requires the underlying system to be
	      fully observed \citep[Sec.~III]{rosas2024}, which is why they
	      secure the observational reading rather than the modal one on
	      their own. The criterion here puts those results to a more explicit philosophical use, as a condition on robust admission
	      that carries explicit intervention indexing, graded verdicts, and
	      downgrade conditions.
	\item
	      Beyond effective realism without downgrade discipline: recent work
	      in ontic structural realism allows regime-dependent, scale-relative
	      ontology at the level of a metaphysical package rather than an
	      individuation criterion \citep{ladymanlorenzetti2024}, but leaves open what
	      operationally distinguishes robust admission from cases that should
	      be downgraded or withheld. Closure supplies that discipline.
\end{enumerate}

With Dennett, the paper
keeps compression realism while refusing to let compression settle
admittance by itself \citep{dennett1991}. With Franklin and Robertson, it
agrees that screening off and novelty are indispensable but argues that
they do not yet separate autonomous from parasitic macro-variables
\citep{franklinrobertson2024}. With Ladyman and Ross, it keeps the
rainforest ambition while insisting that robust admission answer to an
explicit transition criterion \citep{ladyman2007}. And with Ladyman and
Lorenzetti, it takes ontology to be scale-relative and regime-dependent
while remaining answerable to objective modal structure
\citep{ladymanlorenzetti2024}; the present paper then uses interventions
as diagnostics of structural autonomy rather than as metaphysical
primitives. One consequence of intervention indexing
is that the closure criterion can diverge from purely observational
diagnostics: a macro-variable can achieve predictive sufficiency by
tracking an off-partition driver rather than by instantiating autonomous
transition structure
(\S\ref{what-the-formal-machinery-alone-does-not-provide}).

Several neighboring traditions are engaged where they become
dialectically active: Simon and Ando on near-decomposability and
aggregation \citep{simon1962,ando1961}, Wimsatt on non-aggregativity and
reductionist bias \citep{wimsatt2000}, mechanistic accounts of
organized explanation \citep{machamer2000,craver2007}, Rosas et al. on
formal closure diagnostics \citep{rosas2024}, causal-emergence work
\citep{hoel2013,dewhurst2021}, and recent efforts to tighten or reframe
pattern realism and adjacent structuralist debates
\citep{millhouse2022,wallace2024,jiang2024,jiang2025,ladyman2026,meyer2026}.
Each is addressed at the point where its specific contribution bears on
the argument. The narrower claim is that even the strongest of these
views leaves open an admittance question that closure helps make more
explicit and more operational.

\subsection{Intervention, Structure, and the Causation
	Worry}\label{intervention-structure-and-the-causation-worry}

Ladyman and Ross are
skeptical that causation belongs in fundamental metaphysics, and recent
effective ontic structural realism continues that broadly structuralist
posture \citep{ladyman2007,ladymanlorenzetti2024}. By contrast, the
criterion defended here is explicitly intervention-indexed. That can
look like a retreat from structural realism to a Woodwardian
metaphysics of causation, but intervention is not treated here as an
ontological primitive. It serves as a diagnostic probe for structural
autonomy. Admissible
intervention is the most disciplined test of whether a candidate
macro-partition preserves stable relational organization under
perturbation. What the criterion tracks is still structure. The
intervention test asks whether that structure is autonomous at the
candidate grain or whether the apparent autonomy was borrowed from
elsewhere. Woodwardian intervention tests typically operate within an
accepted variable scheme; the criterion here asks the prior question
of when a candidate coarse-graining has earned the right to figure as a
macro-variable at all.

The proposal is compatible with effective ontic structural
realism because Ladyman and Lorenzetti explicitly allow ontology to be
regime-dependent and scale-relative while remaining realist about the
structures discovered there \citep{ladymanlorenzetti2024}. The present
account agrees, but adds a sharper admittance condition: a scale-relative
structure warrants robust ontological standing when its transition
profile is stable under the admissible perturbations that define the
regime. Intervention is used here to supply operational discipline for
effective realism, not to compete with it. On their view ontology is always effective ontology at some
scale and under some modal profile, which does not reduce it to
convenience \citep{ladymanlorenzetti2024}; closure adds a stricter
account of when such scale-relative commitment should be robust,
qualified, or withheld.

What the paper requires is intervention-invariance, as evidence that a
candidate partition tracks stable organization across physically real,
regime-preserving perturbations to boundary conditions, couplings, or
control channels. Interventions supply a methodological route to structural stability. Their use here is not a retreat to fundamental causal ontology.

A more principled objection to intervention-indexing differs from the
familiar complaint that interventions are merely difficult to perform.
Meyer defends the
conditional-probability approach partly on the grounds that it avoids
spatiotemporal presuppositions and preserves metaphysical neutrality,
which matters where the spatiotemporal and causal structure that
interventionist tests presuppose is itself what is in question, as in
quantum gravity \citep{meyer2026}. That argument is a good one, and the reply concedes it rather than rebutting it. Where no admissible intervention class can be specified, the criterion does not deliver a robust verdict; it delivers an indeterminate one, as its downgrade rules require, and that outcome is not a failure of the criterion. The scope
claimed here is induced closure in systems that have a control ecology
(\S\ref{scope-and-non-claims}), and in that scope the observational
reading is not neutral between autonomous and parasitic candidates,
which is what the argument below turns on. Meyer's desideratum and this
criterion are therefore answering different questions: his asks what
travels into domains where intervention talk may not apply, and this one
asks what settles admission where it does. Within that scope, the
neutral reading leaves undecided the autonomous-versus-parasitic cases
that an intervention test can decide.

\subsection{Scope and Non-Claims}\label{scope-and-non-claims}

Closure comes in two forms, stipulated and induced. Stipulated closure
holds when autonomy is fixed by constitutive rules, as in chess or
formal arithmetic; once the rules are set, transitions and consequences
follow by construction. Induced closure holds when autonomy is
discovered through dynamical screening-off in the physical world. The
restriction to induced closure has a principled basis. Induced closure asks whether transition autonomy is sustained by
physical and organizational structure independent of whether anyone
stipulates that it should be; stipulated closure asks only whether a
consequence-relation holds within a rule-system. Institutional kinds
are hybrid, combining constitutive rules with an implementation
substrate, and the criterion tests whether the induced component
sustains the stipulated one under admissible perturbation.
That is why a currency can fail closure while its legal-tender rules
persist unchanged: what degrades is the induced infrastructure of
enforcement and settlement, not the stipulated rule. The scope here is
accordingly restricted to induced closure in spatiotemporal systems,
leaving aside broader debates about abstract objects and a complete
metaphysics of levels; where those debates arise, they are handled only
to protect the central claim from predictable misreadings. Still, the criterion is stated over
state-transition structure: spatiotemporal dynamics are the paradigm
case, not the only intelligible format in which closure questions can be
posed.

The paper offers closure as one part of an account of objecthood, the
dynamical part. It says whether a candidate carries autonomous
transition structure at its grain. Other questions about the same
entities, such as how that structure is produced, what keeps it going,
how we classify it, and how we trace its identity over time, fall to
other accounts, and several current ones fit alongside closure without
competing with it. Contextual emergence treats higher-level
descriptions as licensed by contingent contexts that lower-level
description does not itself supply \citep{bishopatmanspacher2006}; the
regime-indexing defended here is a version of that idea. Process
ontology treats organisms and other entities as stable organizations of
ongoing change \citep{dupre2018,jiang2025}, which is how closure treats
them too: what persists is a transition organization, not a fixed
collection of matter. Structurally oriented hylomorphism fits to the
extent that ``form'' is read as the organizing structure a substrate
realizes \citep{jaworski2016}. Accounts that tie objecthood to physical
binding describe one important way closure is produced in some physical
domains. Functional and historical accounts of artifacts explain why
something is classified as a hammer or a chair \citep{thomasson2007}.
What closure adds to any of these is a check on whether the entity they
describe has dynamics of its own at the relevant grain. If it does,
that autonomy shows up in the transition structure.

The criterion is also not an answer to the special composition
question, which asks when some parts compose a whole
\citep{vaninwagen1990,lewis1986}. It asks when a coarse-graining
carries its own dynamics, and lower-level constituents enter only as
what realizes those dynamics. So the view is compatible with several
positions on composition, including structuralist views on which talk
of objects can be given up in favor of talk of structure
\citep{french2014}. One disagreement does remain. The criterion treats
closed macro-level transition structure as a real feature of the
world. A view that denies this, and holds that only the micro-level
facts are real, disagrees with the paper about substance. A view that
grants it but declines to call the structures ``objects'' disagrees
only about the word, and nothing in the argument turns on the word.

Methodologically, the paper does not depend on the strongest
anti-analytic rhetoric sometimes associated with naturalized
metaphysics. The claim is narrower: closure is a usable philosophical
constraint that can be motivated by dynamical and coarse-graining
practice, whether or not one accepts every meta-metaphysical commitment
in the broader \emph{Every Thing Must Go} (ETMG) program \citep{ladyman2007}.

The paper rejects the view that compression and prediction are
sufficient for realist commitment, but it also rejects a fallback to
utility alone. The claim defended here is closure realism: compression
and projectibility identify serious candidates, while robust commitment
requires transition autonomy under admissible intervention, with regime
and horizon held fixed.
The paper has three argumentative layers: Section 2 states the criterion
and its admissibility discipline; Section 3 gives the exact benchmark in
Markov settings; Section 4 extends the same criterion to non-ideal
settings as a graded protocol. Section 5 addresses the strongest
objections. Section 6 summarizes the gains and limits.

\section{The Criterion: Closure Under Admissible
  Intervention}\label{closure-and-admissibility}

\subsection{Criterion in Plain
	Language}\label{criterion-in-plain-language}

Closure can be stated without heavy formalism. A candidate macro-variable
passes when knowing its current macrostate is enough to determine
its macro-future over a specified horizon and intervention
class. If two microstates inside one macrostate produce different
macro-transitions, closure fails at that grain.

Static patterns fall within the criterion as well. In a static pattern,
micro-fluctuations are still present but screened off, so persistence
is itself a type of macro-transition. The criterion concerns transition
sufficiency rather than descriptive fit. A representation can summarize trajectories elegantly
and still fail to support robust macro-level commitment if it requires
persistent within-class micro-bookkeeping to preserve forecast quality.

Since the formal machinery evaluates partitions of state spaces rather
than physical things directly, an explicit bridge is needed. On the
structural realist commitments defended below
(\S\ref{why-closure-carries-ontological-weight}), realist commitment at
the macro level tracks stable relational and causal organization at the
relevant grain; the partition identifies that organization.

Talk of ``tests'' and ``criteria'' throughout the paper is epistemic
talk, and it should not be read back into the metaphysics. Whether a
partition carries autonomous transition structure under a given regime,
horizon, and intervention class is a fact about the system. It holds or
fails whether or not anyone has checked. The regime is not up to us
either: it is picked out by physical features of the system, such as
timescales, energy scales, and available control channels, that can be
identified without knowing which candidate partition will close
(\S\ref{regime-individuation-and-circularity}). What investigators
contribute is a claim about where such regimes and closed partitions
lie, and that claim can be wrong. The procedural rules below, such as
declaring the regime, horizon, and admissible class before scoring,
govern the claim. They keep the evidence from being tuned to a favored
answer, and they have no role in making the facts. The indexing itself
works like the temperature and pressure conditions on the existence of
ice (\S\ref{regime-dependence-and-projectibility}): it specifies which
modal claim is being made, not whose interests are served. The same
separation appears in the verdict structure of
\S\ref{reporting-and-evidential-discipline}, where a candidate's closure
profile is kept apart from the evidential grade of our access to it.

Closure also does not assign one privileged partition to each physical
system. The same substrate can support several closed coarse-grainings
in overlapping regimes, and each earns standing in its own regime.
Ordinary kind terms need not track these structures one to one. As a
conceptual illustration, suppose a lump of wood and metal formed by
accident and, under the interventions involved in striking and pulling
nails, behaves exactly as a manufactured hammer does. In that regime it
realizes the same transition structure as the manufactured hammer, and
the criterion treats the two alike. Whether the lump is a hammer in the
artifact sense, given that nobody made it for that purpose, is a
separate question. Ordinary language bundles the two questions under
one word, and the criterion answers only the dynamical one.

Several recurring misunderstandings come from running together four
levels.

\begin{enumerate}
	\item
	      World dynamics: the implemented process itself.
	\item
	      Pattern type: stable transition structure supported by that process.
	\item
	      Pattern token: a concrete instance under specific boundary conditions.
	\item
	      Representation: a model, equation, classifier, or narrative that
	      tracks the pattern.
\end{enumerate}

A representation can fail while the pattern type remains real. A token
can fail while the type remains robust. A token can also succeed while
the type is weak, for example in narrow calibration windows. Closure
claims here are primarily type-level claims under specified regimes, then
secondarily claims about token reliability under perturbation.

Closure does not test for every sense in which a pattern might be real.
It tests whether a candidate warrants robust standing under a stated
regime, horizon, and intervention class, and weaker senses of pattern
reality remain available to candidates that do not meet it. Nor does
closure demand perfect transition autonomy, since the graded verdict
structure handles partial cases through qualified and indeterminate
categories. Finally, closure does not replace the special-science kinds
it evaluates. Most of them will earn qualified rather than robust
standing, which is an intended output of the criterion and not a sign
that it has failed.

\subsection{Which Reading of Screening Off Closure
	Secures}\label{why-closure-is-not-screening-off}

A natural rejoinder is that screening off, properly understood, already
excludes parasitic candidates, so closure merely renames an existing
criterion. Much of this can be granted. Closure does not compete with
screening off, and the argument does not require that Franklin and
Robertson overlooked causation or meant something weaker than they said.
Closure supplies a determinate reading of an existing condition,
together with the test that secures it, and adds no further condition.
That counts as a substantive contribution only because the two readings
come apart, and they come apart on candidates that matter.

\begin{enumerate}
	\item
	      Weak reading: observational conditional irrelevance. A variable's
	      value renders further micro-detail probabilistically redundant for
	      the target given the system's actual dynamics. This is a statement
	      about the observational joint distribution, and it is what
	      statistical screening-off diagnostics test.
	\item
	      Strong reading: stability under admissible intervention. The same
	      conditional independence continues to hold under the interventional
	      distribution generated by admissible perturbations that target the
	      candidate partition directly, so that screening off is tracking
	      autonomous structure rather than a shared source.
\end{enumerate}

The claim here is that established screening-off diagnostics generally
secure only the weak reading, while the strong reading is what
warrants robust realist commitment to a macro-variable. Closure operationalizes the
strong reading by adding the admissible intervention class as a
parameter of the test. The two readings agree in the central cases and
diverge when upstream structure lies outside the modeled state space,
or when admissible interventions target the candidate partition
directly while leaving that structure intact. The latent-driver proxy
developed in \S\ref{what-the-formal-machinery-alone-does-not-provide}
is the canonical wedge; the hidden-feedback case
(\S\ref{hidden-feedback-case}) supplies a second wedge whose borrowed
stability has a different source, so the divergence is not tied to one
peculiar configuration.

Accepting the rejoinder therefore means pressing it further. A reading
that no stated test discriminates cannot adjudicate admission, and the
parasitism cases show that the two readings, left undiscriminated,
return opposite verdicts on candidates the debate cares about.
Screening off properly understood is modal screening off, and modal
screening off is transition autonomy. This is one property, described
once in the vocabulary of conditional independence and once in the
vocabulary of dynamics. The concession still leaves the criterion with
work to do. Its contribution is the test and the discipline that goes
with it; it claims no new property. The intervention-indexed reading is
offered as one principled way of making the intended modal content
operationally discriminating, which is why the argument is best understood as an
explication of rainforest admission from within rather than an
objection to it.

\subsection{Formal Statement}\label{formal-statement}

Let micro-process be \(X_t\) and candidate macro-process be
\(Z_t = g(X_t)\). For horizon \(L\), closure asks whether \(Z_t\) is
sufficient for forecasting \(Z_{t+1}^{L}\) under a fixed regime and
admissible intervention class.

The horizon object here is the \(L\)-step path distribution, not only
one-step prediction. One-step tests can be used as diagnostics, but
realist claims are indexed to the declared horizon target.

Informationally, the predictive side asks whether adding \(X_t\) beyond
\(Z_t\) yields material gain for macro-future prediction.
Interventionally, the causal side asks whether macro-transition laws
remain stable across admissible perturbations without needing
within-class micro-identification.

The two tests do not compete. Each reads the same target, transition
autonomy at the macro level.

\subsection{Why This is More Than
	Compression}\label{why-this-is-more-than-compression}

Compression is necessary but not sufficient for robust commitment to a
macro-variable, because a compressed code can hide
transition-relevant heterogeneity and still
score well on narrow tasks. Closure blocks that loophole by asking
whether hidden heterogeneity matters for onward macro-behavior.

Closure therefore functions as an anti-gerrymandering condition,
rewarding coarse-grainings that carry stable transition structure under
declared constraints rather than convenient coding alone.

The account accepts Dennett's claim that compression
success is evidence of objective patterning while denying that
compression success is the full ontological test \citep{dennett1991},
and it accepts Ladyman and Ross's concern that real patterns should be
projectible and structurally disciplined while adding an explicit
transition criterion for adjudicating borderline cases
\citep{ladyman2007}. The aim is to keep those insights while making
their realist force less permissive.

The same condition answers the Haugeland-style pressure point, since
transition autonomy determines when pattern talk has ontological force
rather than merely descriptive convenience \citep{haugeland1998}.

\subsection{Closure and the Primacy of Physics
	Constraint}\label{closure-and-the-primacy-of-physics-constraint}

Ladyman and Ross's Primacy of Physics Constraint is often read
negatively, as a warning against ontologies that float free of physics.
It also contains a second, equally important pressure. Macro-candidates
must be physically compatible, and they must also avoid redundancy
with lower-level description in the realism-relevant sense
\citep{ladyman2007}. The missing question is what that non-redundancy
amounts to once one gives up the idea that only the micro-level is fully
real.

Closure gives that vague requirement operational content. A macro-candidate is non-redundant when
micro-differences within one macroclass do not alter macro-transitions
under the declared regime, horizon, and admissible intervention class.
That is a stronger claim than descriptive usefulness, and a different
claim from mere compression. A macro-description can be shorter, more
tractable, and even explanatorily attractive while still remaining
redundant in the sense that its apparent autonomy depends on continued
micro-level repair.

The latent-driver case shows the difference. A proxy can look
perfectly projectible under passive observation because upstream
structure keeps its trajectories aligned, yet turn out to be
transition-parasitic: perturb the candidate directly with an admissible
intervention, and within-class micro-differences matter again. The
apparent macro-autonomy disappears. Closure therefore sharpens the Primacy of Physics
Constraint's non-redundancy demand by making it a question about
transition structure rather than about descriptive convenience.

Read this way, closure is not an external supplement imposed on rainforest realism. It says more exactly what it takes for a higher-level
pattern to be physically respectable without collapsing back into
lower-level bookkeeping. The discriminator that matters for realism is
causal autonomy under admissible perturbation; microphysical smallness
does not settle the matter by itself.

\subsection{Admissibility and Hierarchical
	Evaluation}\label{admissibility-and-hierarchical-evaluation}

Closure is always relative to an intervention class, and the standard
objection is that this invites circularity, since someone must pick the
class. The reply offered here is procedural and makes no ambitious
metaphysical claim. Any serious
macro-ontology has to answer the same admissibility question, since
any view that evaluates macro-variables has to say which perturbations,
regimes, and horizons are relevant. Closure does not invent that
burden; it makes the corresponding choices public, comparable, and
downgrade-sensitive. Domain interests help determine which control
channels are worth testing, but those choices are constrained by the
system's dynamics, the target regime, and the physically available
modes of perturbation. Objectivity requires that admissibility classes
be declared in advance, cross-compared, and audited. Fragility across
nearby defensible admissibility choices is itself evidence for
downgrade. Robust standing therefore turns on stability of
macro-transition structure across the perturbations the regime can
physically sustain, and fit to observed trajectories is not enough on
its own.

This procedural discipline is the framework's main defense against
gerrymandering. A candidate does not earn robust standing simply
because one carefully chosen intervention class makes it look good.
Human interests do not disappear from the procedure, but the discipline
prevents them from silently protecting a favored partition from
downgrade.

Admissibility is fixed upstream of commitment verdicts by three
constraints:

\begin{enumerate}
	\item
	      Epistemic admissibility: interventions are measurable and
	      implementable by bounded agents.
	\item
	      Dynamical admissibility: interventions preserve the target regime
	      class.
	\item
	      Explanatory admissibility: interventions target variables with
	      cross-context generalizability and non-trivial counterfactual reach,
	      rather than one-off manipulations tuned to a single episode.
\end{enumerate}

These constraints are physically anchored in available control channels
and implementation structure, not in analyst preferences about favored
ontologies. The explanatory constraint is set before closure verdicts and
does not assume in advance which candidate partition will prove stable.
This makes the interventionist notion of a ``possible experiment''
concrete and bounded rather than leaving it at conceptual possibility
\citep{woodward2015}, and aligns with \citet{woodward2021}'s
proportionality demand on variable choice. Explanatory admissibility is
stronger when its control channels already exist in the domain and
replicate across sites, weaker when success depends on narrow one-off
tuning. Domain expertise shapes regime identification and control-channel
selection, but that shaping is declared, auditable, and separated from
partition scoring.

Pearl's $do(x)$ surgery is defined mathematically regardless of
physical feasibility, which is productive for causal inference
\citep{pearl2000}. Robust realist commitment requires tighter
constraints. A partition that closes only under interventions no bounded agent could
implement without destroying the target regime does not support the
intervention-guiding structure that warrants realist commitment. The
restriction concerns which of Pearl's outputs carry ontological
weight, not the framework itself.

Nor is the admissible class arbitrary. A system's boundaries are probed constantly by
naturally occurring perturbations (thermal noise, diffusion gradients,
mechanical impacts), and a cell membrane screens molecular fluctuations
whether or not an experimenter is watching.
These perturbations are exogenous to the candidate partition even when
internal to the larger system, so the test inherits the exogeneity that
standard interventionist accounts demand without requiring human agency
to secure it.

Evaluation order is fixed by the canonical selection discipline
(\S\ref{procedural-safeguards-in-non-ideal-cases}): intervention
classes are fixed from independent physical constraints and established
control practice before closure is scored, and post hoc revisions
restart the evaluation.

Taken together, these requirements give a short acceptability test. An
admissibility class is acceptable only when:

\begin{enumerate}
	\item
	      interventions use physically available control channels;
	\item
	      interventions preserve the regime being evaluated;
	\item
	      no intervention is one-off or directly partition-protective;
	\item
	      verdicts are checked across nearby defensible alternatives.
\end{enumerate}

Classes that fail any of these conditions cannot carry the evidential
weight of a robust verdict.

\subsection{Regime Individuation and
	Circularity}\label{regime-individuation-and-circularity}

If regime individuation already commits to which macro-variables matter,
the protocol risks circularity: closure testing presupposes regime
identification, but regime identification might presuppose partition
outcomes. Part of the reply is structural. The criterion is
designed to sort cases into three classes rather than to force robust
verdicts everywhere. In the first class, regime specification is
uncontroversial across candidate partitions, and the criterion
adjudicates robustly. In the second, regime specification is contested
but iteration stabilizes under the selection discipline
(\S\ref{procedural-safeguards-in-non-ideal-cases}), and reflective
equilibrium converges on a verdict. In the third, regime specification
is contested and iteration does not stabilize, so the criterion issues
an indeterminate verdict. For a case of that kind, indeterminacy is
the correct output and not a failure of the criterion.

The other part of the reply is procedural. Defensible regime
specifications draw on quantities characterizable without presupposing
which candidate partition will succeed. Temporal parameters, energy
scales, flow rates, and enforcement intensities are typically
identifiable through physical indicators that competing candidate
partitions can agree on before either is evaluated. A running example,
used again at two later points, illustrates this. Consider two rival partitions of weekday commuter traffic in a
metropolitan corridor: \(P_{lane}\), which tracks density and flow by
lane segment, and \(P_{tod}\), which tracks flow volume in clock-time
buckets calibrated to typical commuter patterns. Both are choices a real
modeler might defend, and their in-sample one-step fit can be similar.
The regime is characterized by clock schedules, total vehicle counts, and
road network topology, all quantities that investigators favoring either
partition can agree on before either closure test is run. Regime
specifications of this kind are anchored in
lower-level or cross-cutting physical indicators rather than in the
outcome of the very test they enable. Regime
specifications should be grain-neutral across competing candidate
partitions: a regime can itself be a macro-level description while
still not presupposing which finer-grained partition will close
within it.

Independence is imperfect in some domains, particularly biology and
social science, where macro-variable choice and regime characterization
can be entangled (psychiatric diagnostic categories, species concepts:
the variables that pick out the regime are themselves candidates whose
macro-autonomy is under test). The criterion is built to handle such
entanglement, not undone by it. When alternative, independently
motivated regime specifications produce divergent closure verdicts, the
right result \emph{is} downgrade to at most qualified, because the
evidential situation does not yet support a stronger commitment. A
candidate whose apparent autonomy is parasitic on the choice of regime
characterization should not receive robust standing.

Consider a contested case as a stress test, sketched schematically and
not settled here. Suppose major depressive disorder is approached under
three independently motivated regime characterizations: a symptom-cluster
regime tied to diagnostic instruments, where intervention evidence on
pharmacological and psychotherapeutic levers might support partial
closure within a bounded clinical horizon; a neurobiological regime
indexed to specific circuit dysfunctions, where the same label may
fragment across heterogeneous mechanisms; and a social-functional regime
indexed to role-impairment trajectories, tracking a third transition
profile. If the case has that shape, the criterion returns qualified
standing, with explicit downgrade for any commitment that depends on
which characterization is accepted upstream. Whether it has that shape is
a question for philosophy of psychiatry, and nothing here settles it.
The example offers no verdict on this kind. It shows only what the
criterion does with a kind whose regime is contested: it declines to
launder that kind into robust standing, which is the response the
criterion is designed to give and should not be read as evasion.

Regime individuation and partition evaluation are iterative rather
than viciously circular: initial regime characterizations are
provisional, and closure testing can refine them.
The process stabilizes into reflective equilibrium in Goodman's sense
\citep{goodman1955}: if iteration does not stabilize across nearby
defensible starting specifications, the result should be treated as
indeterminate rather than forced into a robust verdict. Here
stabilization means that comparative ordering across independently
specified candidate partitions stops shifting under small, defensible
protocol variations; it does not presuppose one privileged grain in
advance. Iteration is governed by the same selection discipline
(\S\ref{procedural-safeguards-in-non-ideal-cases}) as initial scoring:
each revision restarts the evaluation rather than silently inheriting
prior results. The framework does not guarantee a unique answer in
every domain. In hard cases it sometimes relocates disagreement without
resolving it, and that is the intended outcome. What it aims to do is distinguish
cases where robust commitment is warranted from cases where the
evidential situation supports only qualified or indeterminate standing.

A more global form of the worry asks what selects the right regimes and
contexts in the first place. Some sufficiently narrow regime can
probably be found in which almost any candidate looks stable, so it
might seem that closure needs a privileged scientific task to tell it
which contexts count. It does not, because the protocol checks two
kinds of robustness that do different jobs.

Robustness across the admissible intervention class is part of what
closure is. The criterion quantifies over every within-class
realization that an admissible intervention can reach
(\S\ref{leakiness-as-canonical-target}). A candidate that holds up
under some admissible perturbations and fails under others is therefore
not closed under that class, and this failure is a fact about the
candidate rather than a gap in the evidence.

Robustness across nearby regimes, horizons, and modelling choices plays
a different role. It is not part of closure in a single regime, since a
candidate can be closed in one regime and leak in the next, as ice does
across a phase boundary. It is instead evidence that the regime under test is a genuine
feature of the system's dynamics rather than one drawn to suit a
favored partition. A regime cut so narrowly that a favored candidate
looks stable within it, and nowhere nearby, will usually fail
projectibility under small shifts in context
(\S\ref{regime-dependence-and-projectibility}), and a regime whose
boundaries were drawn by looking at closure scores violates the
selection discipline outright. A closure verdict is therefore always
relative to a regime, but the regime has to be identified from the system's physical features and
has to hold up under nearby variation. No further task-level
justification of the context is required, because the contexts that
matter are the ones in which the system actually supports autonomous
transition structure.

``Regime dissolution,'' the event that marks the outer edge of a
regime, is not partition-relative. It can be characterized by
cross-cutting physical indicators that do not presuppose the winning
partition: energy-scale crossings, phase boundaries, control-response
failure at established channels, and breakdowns in the separation of
timescales that defined the regime. These indicators register whether
the background dynamics have changed, not whether a favored partition
continues to score well.

\subsection{Admissibility Disputes}\label{admissibility-disputes}

Admissibility disputes are expected, especially across domains. The
adjudication rule is to compare candidate intervention classes by what
they physically permit and whether they preserve the same target regime.
Boundary-pressure perturbations can be admissible for storm dynamics,
while molecule-by-molecule remote rewriting is not. In institutional
settings, changing enforcement intensity can be admissible, while
instant arbitrary rewriting of all agent commitments is not. The choice
has ontological consequences. A storm satisfies closure under
boundary perturbations because its pressure organization and rotational
structure determine macro-transitions without tracking individual
molecules. Under molecule-by-molecule rewriting, the same candidate
fails, because admissibility now permits interventions that reach inside
the macroclass and exploit within-class differences. Between the two
cases only the admissible control channel differs, and the verdict
follows it. For this reason admissibility must be fixed before any
candidate is scored (\S\ref{admissibility-and-hierarchical-evaluation}).

The issue is distinct from regime individuation. Regime individuation
asks which background dynamics define the target setting; admissibility
individuation asks which intervention classes count as ontologically
relevant probes within that setting. Explanatory admissibility is tied to available
control channels and implementation structure, not to interventions that
happen to vindicate a favored partition, and this answers the
circularity worry here as well.

Front-end specification of regime, target, and admissible control
channels is often interest-guided in a perfectly ordinary way. That does
not make closure verdicts interest-relative once those specifications
are fixed and openly defended. A short template grid (boundary nudges,
coarse actuator controls, policy levers) can make admissibility
comparison public and repeatable in domains that lack an established
control lexicon.

Interventions that directly target preservation of partition labels,
rather than system variables, are inadmissible without exception,
because they trivialize closure by design. The test asks whether
transition structure is autonomous under physically meaningful controls,
and protecting labels by stipulation says nothing about that.

When two intervention classes are both admissible and target the same
regime, the criterion permits plural testing. Verdicts should report
which class was used and how sensitive results are across nearby
defensible classes; sensitivity to a narrowly tuned class is grounds for
downgrade. In the commuter corridor of
\S\ref{regime-individuation-and-circularity}, ramp metering and
speed-limit controls count as admissible because they are physically
available control channels in the regime. Vehicle-by-vehicle remote
rewriting is excluded because it is not part of the regime's actual
control ecology; that it would embarrass \(P_{lane}\) plays no role in
the exclusion.

\subsection{Predictive and Interventional
	Closure}\label{predictive-and-interventional-closure}

Predictive and interventional closure should be distinguished but not
separated. Predictive closure concerns whether macrostate information
screens off transition-relevant micro-detail for macro-future
forecasting; interventional closure concerns whether macro-transition
structure remains stable when the system is probed.

The evidential relation between them is asymmetric. Strong
interventional closure typically implies predictive adequacy for the same
target and horizon, while predictive adequacy alone can persist in cases
where interventional stability fails under distribution shift. For that
reason, the criterion treats predictive evidence as important but not
final.

The asymmetry also sets the burden of proof. A robust verdict needs
either direct interventional support or compelling indirect evidence
that tracks intervention-relevant invariance. Otherwise, the responsible
verdict is qualified or indeterminate.

Indirect evidence can still be substantial: natural experiments,
intervention-like policy shifts, quasi-experimental contrasts, and
systematic failure under exogenous disruption can all reveal whether a
candidate partition preserves its own transition structure across
admissible changes. These sources do not collapse the distinction
between predictive and interventional closure. They are useful because
they bear on intervention-relevant invariance and not on in-sample fit
alone.

In biology, economics, and the social sciences, direct manipulation is
often sparse, ethically constrained, or expensive. The ontological
target nevertheless remains interventional closure, so where the
evidence for that target is thin, robust verdicts should be rare. The
graded verdict structure is designed to yield that result; it does not
relax the target to observational fit when interventional evidence is
scarce.

\subsection{Why Closure Carries Ontological
	Weight}\label{why-closure-carries-ontological-weight}

A reviewer can still ask why transition autonomy should count as robust
macro-level standing rather than as a mere success condition for
modeling. The reply rests on two of the paper's explicit commitments.
Under structural realism, what ontology should track is stable
relational and causal organization rather than intrinsic micro-identity
as such. Under interventionism, what counts as causally relevant
structure is structure that supports stable manipulation and control.

This fits the regime-dependent, scale-relative orientation defended by
effective ontic structural realism. The proposal seeks a sharper
criterion of when a scale-relative structure has earned commitment, and
when it should be downgraded, within a broadly structural realist
framework \citep{ladymanlorenzetti2024}. It does not ask for a return to
one privileged level or to thick intrinsic objects.

The meta-philosophical stance here follows what Woodward calls the
methodological path to ontology: causal-interventionist tools are used
to determine what earns realist commitment, rather than seeking
independent metaphysical ``truth-makers'' or ``grounds'' that stand
apart from the practices of manipulation and prediction \citep{woodward2015}. On that path, scientific ontology disciplined by
closure is the only ontology required for commitment verdicts at this
grain. This is a consequence of the stated commitments rather than an
independent argument against those who reject them. A grounding
theorist who holds that stable dynamical autonomy under admissible
intervention still stands in need of a further metaphysical foundation
is not making a mistake by the lights of their own view; they are
declining the methodological path at the outset, so the disagreement
concerns that path rather than closure. Readers who
take it will find that closure settles what commitment requires;
readers who do not can still take the anti-gerrymandering result of
\S\ref{what-the-formal-machinery-alone-does-not-provide}, which does
not depend on the ontological reading.

A persistent worry is that prediction and interventional control are
merely epistemic lenses, leaving the underlying ontology untouched.
This confuses the test with what is being tested. Prediction and
intervention are diagnostic tools that reveal pre-existing statistical
structure. The conditional independencies that allow a macro-state to
screen off micro-details exist whether or not anyone computes them.
Locality, finite signal propagation, and thermodynamic limits create
these independencies as physical facts. Our predictive models succeed
when they align with this objective causal topology, and fail when they
do not. Closure is therefore an ontological criterion: it identifies
where the world's phase space objectively decouples, which need not
coincide with where our minds prefer to chunk it. Under structural
realism, such decouplings count as ontological because they are the
stable relational structure that realist commitment was meant to track
in the first place, and not because they license any particular
representation. A criterion that identifies them picks out what the
ontology at that grain consists in, rather than a consequence of it.

Given those commitments, closure is the criterion that identifies when a
candidate macro-description tracks stable structure at the level where
explanation and intervention are being assessed. It is not a convenient
heuristic layered on top of ontology. This is how the view remains
compatible with microphysical completeness while still defending
non-redundant macro commitment. Nor does the paper smuggle ontology out
of modeling success. It defends a specific thesis about which kind of
modeling success is ontology-relevant, namely transition autonomy under
admissible perturbation.

What the criterion measures is a feature of the candidate rather than a
label imposed from outside the formalism. A candidate's closure profile,
how much transition autonomy it has under the declared constraints, does
not depend on anyone's confidence in it. The verdicts built on that
profile are a further matter, since they answer to the state of the
evidence as well (\S\ref{reporting-and-evidential-discipline}); the
ontological claim here attaches to the profile.

\subsection{Closure and
	Underdetermination}\label{closure-underdetermination-and-objecthood-discipline}

One remaining concern is underdetermination. Even after the criterion is
stated, a reviewer may argue that many distinct coarse-grainings can be
tuned to look acceptable on available data. If so, closure might seem to
collapse back into model choice rather than commitment discipline.

Three questions need to be kept apart.

\begin{enumerate}
	\item
	      Which candidate partitions are descriptively serviceable in a dataset?
	\item
	      Which candidates remain transition-autonomous under admissible
	      perturbation?
	\item
	      Which of those candidates remain stable under modest horizon and
	      admissibility variation?
\end{enumerate}

Underdetermination is strongest at the first question and weaker at the
second. It is often weakest at the third. Many candidates can fit
observed trajectories, far fewer preserve counterfactual structure when
the regime is probed, and fewer still remain stable across nearby
defensible setups.

The account does not claim that every domain yields one uniquely
privileged partition. It claims that realist commitments should be
constrained by transition autonomy and robustness, rather than by fit
alone. When underdetermination persists after those filters, the
responsible result is qualified or plural commitment under transparent
constraints, not forced monism and not unconstrained relativism.

When exact closure is unavailable, leakiness-centered diagnostics
estimate comparative distance from closure, and commitment strength
tracks the stability of that evidence.

\section{Exact Benchmark: Strong
  Lumpability}\label{exact-benchmark-strong-lumpability}

From Rosas et al., the paper takes the formal diagnostics of
informational, causal, and computational closure, and the
strong-lumpability benchmark \citep{rosas2024}. It adds an
intervention-indexed criterion for when such diagnostics warrant robust
commitment to a macro-variable, together with downgrade rules for
non-ideal cases. The formal results are read as the ideal limit of the
criterion rather than as the criterion itself.

Strong lumpability serves as a benchmark: a clean exact case in Markov
settings where closure can be stated and checked without ambiguity.

Let partition cells under \(g\) be macroclasses. Strong lumpability
holds when microstates within the same macroclass induce identical
transition probabilities to all macroclasses. This is the standard
strong-lumpability condition; cf. the presentation attributed here to
Kemeny and Snell \citep{kemeny1960}. When this condition holds,
macro-transitions are autonomous by construction.

More precisely: for any two microstates \(x, x'\) in macroclass \(C_i\)
and any macroclass \(C_j\), strong lumpability requires
\(\sum_{y \in C_j} P(y \mid x) = \sum_{y \in C_j} P(y \mid x')\). The
induced macro-kernel is then well-defined independently of
representative choice, so the macro-process is itself Markov. A
partition that fails this condition mixes transition-heterogeneous
microstates and therefore lacks macro autonomy.

The Markov condition does not assume history is irrelevant; it assumes
only that the relevant state description has been fixed. In
memory-bearing systems, that description may need to be history-enriched
before the closure question is asked. That does not trivialize
non-Markovian cases or the formal difficulty of establishing
lumpability in them; it states the standard point that the closure
test applies to the least state description that properly bounds the
relevant memory.

Computational mechanics puts the same point in sharper language. For
any macro-process \(Z\), two prediction machines can be defined. The \(\varepsilon\)-machine is the minimal
model that predicts \(Z\)'s future from \(Z\)'s own past, using only
macro-level information \citep{shalizi2001}. The \(\upsilon\)-machine is the minimal model
that predicts \(Z\)'s future from the full micro-past \(X\), using
everything available \citep{rosas2024}.
Closure holds when these two machines are equivalent: the
\(\varepsilon\)-machine already captures everything the
\(\upsilon\)-machine knows about macro-futures, and extra
micro-information is redundant for the macro target. This gives an exact
ideal where closure is not approximate. The same formal point is developed
in Rosas et al.'s treatment, where strong lumpability corresponds to the
exact case in which the macro-level causal-state machine coincides with
the micro-conditioned prediction machine, so that no residual
micro-information adds predictive leverage
\citep[Thm.~3 and Prop.~4, pp.~27--28]{rosas2024}.

Exact strong lumpability is uncommon in open, noisy, and path-dependent
systems. That rarity does not undermine the criterion, which uses the
exact case as its ideal limit.

\subsection{Minimal Formal
	Illustration}\label{minimal-formal-illustration}

The anti-gerrymandering role can be shown with a minimal symbolic
example. Let microstates be \(\{x_1,x_2,x_3,x_4\}\) with transition
rows:

\[
	\begin{aligned}
		P(x_1,\cdot) & = (\alpha,\beta,0,0),  \\
		P(x_2,\cdot) & = (\alpha,\beta,0,0),  \\
		P(x_3,\cdot) & = (0,0,\gamma,\delta), \\
		P(x_4,\cdot) & = (0,0,\gamma,\delta),
	\end{aligned}
\]

with \(\alpha+\beta=1\) and \(\gamma+\delta=1\).

Partition \(A\) groups \(\{x_1,x_2\}\) and \(\{x_3,x_4\}\). Partition
\(B\) groups \(\{x_1,x_3\}\) and \(\{x_2,x_4\}\). In \(A\), members of
each macroclass have matching onward profiles to macroclasses, so
macro-transitions are autonomous. In \(B\), members of one macroclass
differ in onward profiles whenever \(\alpha \neq \gamma\), so the
macro-label hides a transition-relevant difference.

Partition \(B\) can still be made predictive by adding repeated
within-class bookkeeping. A partition is \emph{high-maintenance} when
reliable macro-prediction repeatedly forces refinement of within-class
distinctions, so that yesterday's macrostate labels must be supplemented
by new micro-bookkeeping to retain accuracy. The criterion excludes such
partitions from robust standing.
On a fixed sample trajectory, this bookkeeping can make Partition \(B\)
look nearly as predictive as Partition \(A\). The difference appears
under admissible perturbation: \(A\) preserves macro-transition structure
without renewed microstate tracking, whereas \(B\) does not.

Both partitions are definable, but only \(A\) preserves transition
autonomy, so closure separates legitimate macro-candidates from
aggregates that are merely codable.

\subsection{How the Formal Machinery is Used
	Here}\label{what-the-formal-machinery-alone-does-not-provide}

The mathematical property and the philosophical criterion are different things. The criterion does work that the formal framework does not attempt and, for its own purposes, does not need to attempt. The central addition is intervention indexing.
Causal-state constructions classify
histories by equivalence of future distributions under the observed
process, which is primarily an observational-predictive equivalence
relation \citep{shalizi2001}. The criterion here adds an
admissible intervention class as a parameter of the closure test.
Realist claims are indexed to what the system does under
perturbation, not only under passive observation, which separates
the criterion from a purely statistical diagnostic.
Rosas et al.'s framework distinguishes informational, causal, and
computational closure as complementary formal diagnostics
\citep{rosas2024}, and their Theorem 1 establishes that for spatial
coarse-grainings, under the assumptions specified there, informational
and causal closure coincide \citep[Thm.~1]{rosas2024}. In that
framework \emph{causal closure} is defined as equivalence between the
coarse-grainings induced by a process's \(\epsilon\)-machine and its
\(\upsilon\)-machine \citep[Def.~3]{rosas2024}, and \emph{information
closure} as the vanishing of \(I(\vec{X}_t; \vec{Z}^L_{t+1} \mid
\vec{Z}_t)\) for all \(L\) \citep[Def.~4]{rosas2024}. Both are defined
over the observed process. Rosas et al. are themselves explicit about
what this does and does not license: causal closure as defined there
provides counterfactual guarantees only if the underlying system is
fully observed, and otherwise ``ought to be understood in the Granger
sense, i.e.\ as a best guess given the available knowledge''
\citep[Sec.~III]{rosas2024}. The equivalence the theorem establishes
is therefore one between two readings of observational structure, not
between observational structure and interventional invariance under
perturbation external to the modeled process. The criterion here is
downstream of that result, and it takes their caveat as its point of
departure. It treats those observational diagnostics as strong evidence
for closure, but ties realist commitment to a further
admissible-intervention stability condition and to the absence of
hidden micro-identity bookkeeping. The two diagnostics agree in the
central cases and can diverge when relevant upstream structure lies
outside the modeled state space, or when admissible interventions
target the candidate partition directly while leaving such structure
intact. Such divergence does not count against their result, because it
falls in the region their own full-observation proviso exempts. Where
the structure sustaining apparent macro-autonomy lies off-partition, the
Granger reading is the operative one, and the counterfactual guarantee
lapses at exactly the point where parasitism becomes possible. The
latent-driver case developed below is the first instance; the
hidden-feedback case (\S\ref{hidden-feedback-case}) shows the same
failure signature from a structurally different source.

The graded assessments the criterion issues come in two parts. A
candidate's closure profile reports where it sits among more and less
stable closure relations, and that is a claim about the candidate. The
commitment verdict built on the profile also answers to the evidence,
so a downgrade can reflect either a fragile candidate or thin access to
it (\S\ref{reporting-and-evidential-discipline}).

\subsubsection{A Latent-Driver Proxy}\label{latent-driver-proxy}

The first case is the latent-driver proxy, where observational
screening off cannot discriminate autonomous macro-structure from
borrowed macro-structure. Consider a composite variable over subsystems
that are pinned by an upstream driver process the partition does not
track. The driver need not be a single cause that synchronizes the
subsystems; that is only one way the structure can be realized. What
matters is that the driver process is upstream of both the composite's
present state and its future state. The apparent dependence of the macro-variable
on its own past is then underwritten by off-partition structure. Under
passive observation, the composite's macro-state predicts its own
future as well as the full micro-state does, because the driver pins
the subsystems and thereby makes the proxy's present state carry
exactly the upstream information its future turns on. An
\(\varepsilon\)-machine built on the composite finds no residual
micro-information worth adding, and screening-off diagnostics register
a clean pass. Can the screening-off-plus-novelty package handle this
case?

Franklin and Robertson's criterion asks whether the macro-variable
enters into a genuine macrodependency that renders some lower-level
detail conditionally irrelevant while contributing something novel at
the macro level \citep[\S3]{franklinrobertson2024}. Their screening-off
condition has two parts. The dependence part, \emph{unconditional
relevance}, requires \(P(A \mid LLD) > P(A)\): the lower-level
description is relevant to the macro-fact. The independence part,
\emph{conditional irrelevance}, requires in its strict form that
\(P(A \mid B \,\&\, LLD) = P(A \mid B)\), so that all the influence of
the micro-detail is mediated through the macrostate
\citep[\S3.1]{franklinrobertson2024}. Their novelty condition is
more demanding than it may first appear, and a weaker reading would
make the case below too easy. Novelty requires \emph{distinct}
macrodependencies: a binary star
system screens off the individual stellar trajectories, yet they deny
it admission, because the dependencies describing it ``are just the
same nomic dependencies that describe the trajectories of the two
individual stars'' \citep[\S3.2]{franklinrobertson2024}. A dependence
that is merely statable at the macro grain therefore does not satisfy
novelty. Novelty requires a macrodependency of a different form from
the micro-dependencies that realize it, and its mark is
cross-classification, where divisions at the macro level cut across
divisions below.

That is a real constraint on the pressure case, and it excludes the
simplest candidates. A macro-variable that is a relabeling of a single
driver inherits that driver's law and so fails novelty for the same
reason the binary star does. The proxy below is therefore more
elaborate than the simplest candidates: it is the simplest one that
survives their whole package.

Consider an upstream driver process with two components, \(U^1\) and
\(U^2\), independent of each other, each a symmetric two-state Markov
chain on \(\{0,1\}\) that flips with probability \(q\), where
\(0 < q < 1/2\). Two subsystems \(A\) and \(B\) track their respective
components. The system's specified dynamical ecology includes one
control channel, a reset line on \(B\) with indicator \(R_t\) and value
\(C_t\). The update is simultaneous, and every variable at \(t+1\) is
fixed by variables at \(t+1\) or \(t\):
\[
	U^i_{t+1} = U^i_t \oplus N^i_{t+1}, \qquad
	A_{t+1} = U^1_{t+1}, \qquad
	B_{t+1} =
	\begin{cases}
		C_{t+1}   & \text{if } R_{t+1} = 1, \\
		U^2_{t+1} & \text{if } R_{t+1} = 0,
	\end{cases}
\]
where \(N^1, N^2\) are independent Bernoulli\((q)\) noise terms. The
micro-state is \(X_t = (U^1_t, U^2_t, A_t, B_t)\), and the reset line is
an exogenous control input rather than a component of \(X_t\): the
controlled process is Markov on \(X_t\) once \((R_t, C_t)\) is fixed.
The candidate macro-variable is the agreement of the two subsystems,
\(Z_t = A_t \oplus B_t\), taking the value \(0\) when they agree and
\(1\) when they disagree.

The reset channel is part of the system's specification rather than an
intervention imagined from outside it. In the observed process
\(R_t = 0\) for all \(t\), so the subsystems are pinned to their
components, \(A_t = U^1_t\) and \(B_t = U^2_t\), and
\(Z_t = U^1_t \oplus U^2_t\). The restriction of the observed process to
the synchronized manifold is therefore a fact about which channels
happen to be exercised, not about what the dynamics permit.

Write \(V_t = U^1_t \oplus U^2_t\) for the driver parity and
\(r = q^2 + (1-q)^2\) for the probability that the two noise terms
agree. Substituting the update rule gives the identity the whole case
turns on: with \(R_{t+1} = 0\),
\[
	Z_{t+1} = V_t \oplus (N^1_{t+1} \oplus N^2_{t+1}),
	\qquad\text{so}\qquad
	P(Z_{t+1} = V_t) = r .
\]
The macro-variable's future is fixed by the \emph{driver} parity, not by
its own current value. The two coincide exactly when the channel is
dormant, since \(Z_t = V_t\) whenever \(R_t = 0\). Everything below
follows from that gap.

\begin{enumerate}
	\item
	      Unconditional relevance, taken with their indexing. Their
	      condition ranges over the lower-level descriptions that subvene
	      a given macrostate \(B\) and influence a macro-fact \(A\), which
	      together form an equivalence class
	      \citep[\S3.1]{franklinrobertson2024}. Fix \(A\) as
	      \(Z_{t+1} = 0\) and \(B\) as \(Z_t = 0\). In the observed
	      process the descriptions subvening \(B\) are exactly those with
	      \(V_t = 0\), namely \((0,0)\) and \((1,1)\), and each raises the
	      probability of \(A\) from its unconditional value of \(1/2\) to
	      \(r\), since \(r - 1/2 = 2(q - 1/2)^2 > 0\) for
	      \(q \neq 1/2\). So \(P(A \mid LLD) > P(A)\) holds across that
	      class, which is what their condition asks. Descriptions with
	      \(V_t = 1\) lower the probability of \(A\) to \(1 - r\), but
	      they subvene \(Z_t = 1\) and so belong to the class for a
	      different macrostate.
	\item
	      Conditional irrelevance, in the strict rather than the
	      approximate form. When \(Z_t = 0\), the only micro-states the
	      observed process ever visits are those with \(V_t = 0\), and
	      both yield \(P(Z_{t+1} = 0) = r\). When \(Z_t = 1\), the only
	      visited realizations have \(V_t = 1\), and both yield
	      \(1 - r\). Hence
	      \(P(Z_{t+1} \mid Z_t, X_t) = P(Z_{t+1} \mid Z_t)\) exactly.
	      The unvisited realizations do not disturb this verdict. The
	      realizations of \(Z_t = 0\) with \(V_t = 1\)
	      receive probability zero in the observed process, so the
	      conditional probability is undefined on them and the strict
	      condition holds over them only vacuously. The condition is
	      silent about these heterogeneous realizations, and they are
	      exactly the ones on which the closure question turns. So the
	      proxy satisfies the strict condition on every realization for
	      which the condition says anything at all, rather than
	      approximating screening off in the way a deviant microstate
	      would.
	\item
	      Novelty, structurally rather than numerically. \(Z\)'s
	      macrodependency relates parity classes to later parity classes,
	      and neither its relata nor its transition partition coincide
	      with those of either component process. The macrostate
	      \(Z = 0\) lumps \((0,0)\) with \((1,1)\), which cuts across the
	      partition induced by either component's state, so the
	      equivalence classes over which \(Z\)'s transitions are defined
	      are not the classes over which \(U^1\)'s or \(U^2\)'s
	      transitions are defined. This is cross-classification in
	      Franklin and Robertson's sense, and it is why the case is not
	      the binary star, whose centre of mass obeys the very
	      dependencies its components obey. The numerical difference is a
	      symptom of this rather than the argument for it: \(Z\) is a
	      symmetric chain with persistence \(r\) against each component's
	      \(1 - q\), and \(r - (1-q) = q(2q-1) \neq 0\) for
	      \(q \neq 1/2\). Were the parameters to coincide, the
	      macrodependency would still not be either component's, because
	      it would still be defined over different relata.
\end{enumerate}

A defender may resist at the third condition, holding that parity is a
mere algebraic consequence of two independent processes and so lacks the
relevant kind of novelty. That reading may be defensible, but it is not
the published one. If novelty
requires more than cross-classification and non-identity of dependency
profiles, whether organizational integration, explanatory unity, or
something else, that further requirement has to be stated independently,
because the probabilistic criterion as given does not supply it. A
defender who adds it concedes the paper's point: the stated test
underdetermines the verdict, and the verdict is settled by a commitment
the test does not contain.

The stated implementation of the criterion therefore returns a pass.
This is weaker than saying the proxy belongs in the rainforest, and the
weaker claim is the one the argument needs: the conditions as stated,
evaluated on the observed distribution, do not exclude this candidate. Whether Franklin and Robertson's broader modal
intentions would exclude it is exactly the question at issue, and the
thesis here is underdetermination rather than refutation.

Now consider the micro-realizations of \(Z_t = 0\) that the observed
process never visits: those with \(V_t = 1\). Take drivers at
\(U^1_t = 0\) and \(U^2_t = 1\), and actuate the reset line at \(t\)
with \(R_t = 1\) and \(C_t = 0\). Then \(B_t = 0\), \(A_t = U^1_t = 0\),
and \(Z_t = 0\), while the drivers are untouched and \(V_t = 1\). From
there \(P(Z_{t+1} = 0) = 1 - r\), against \(r\) for every observed
realization of \(Z_t = 0\), and \(r \neq 1 - r\) for \(q \neq 1/2\). Two
micro-realizations of one macrostate carry different macro-futures, and
the update rule shows why: the future tracks \(V_t\), while the
macrostate tracks \(V_t\) only for as long as the channel stays dormant.

Building the channel into the specification forecloses the natural
objection that the perturbation forces the system off the manifold on
which the macro-variable was defined, and so abandons the regime rather
than testing it. Actuating the reset line is an ordinary operation of
the system's own control ecology. The dynamics are unchanged throughout,
the channel falls dormant again immediately, and the subsystems are
repinned to their components at the next step, so the deviation is a
one-step transient. The synchronized manifold was never a constraint the
dynamics imposed. It was a restriction on the observed distribution,
produced by leaving an available channel unused, and the criterion is
designed to detect exactly that difference.

The countermodel invites a harder version of the objection: any
macro-variable would fail closure if one were free to posit enough
internal reset channels. That is true, and it is the reason for the
admissibility constraints of
\S\ref{admissibility-and-hierarchical-evaluation}. A
channel cannot be introduced after the verdict is seen, must belong to
the regime's independently specified control ecology, must preserve the
target regime, and must survive comparison with nearby defensible
intervention classes. A reset line invented to embarrass a candidate
fails the first and second of these immediately.

A second worry concerns ordinary dependence. Organisms, artifacts, and
institutions all depend on their surroundings: on gravity, on a
temperature range, on nutrient flows, on the hand that swings the
hammer, on supporting practices. If any dependence on structure outside
the partition counted as parasitism, the wedge would condemn nearly
every special-science kind. Parasitism in the present sense is much
narrower than dependence, and the two ways an environment can support a
candidate show where the line falls.

Some environmental conditions are held fixed by the regime: gravity,
the temperature range, the institutional background. These are part of
the regime's specification. The candidate's closure is assessed against
them, and a change in them is a change of regime rather than a
perturbation within it. Other influences vary and act on the candidate
as inputs, as the swing acts on the hammer or a nutrient supply on a
cell. These are treated as the reset line is treated in the model
above: as exogenous inputs that the transition law takes as given, not
as part of the candidate's state. The closure question is then whether
the candidate's own macrostate, together with those inputs, fixes its
macro-future. For the hammer and the cell it plausibly does. Given the
swing, the hammer's position and momentum settle what happens to the
nail; the lump's microstructure adds nothing. Dependence of this kind
leaves the candidate's own state doing the predictive work.

The proxy is different. Its drivers are neither fixed background nor
inputs acting on a state that carries its own memory. Given the drivers,
the proxy's present value adds nothing to the forecast of its future,
which is fixed entirely by the driver parity. Its present value looks
predictive only because, in the observed process, it happens to
coincide with that parity. So the proxy's apparent macro-dynamics are
the driver's dynamics seen through a correlate, and an admissible
intervention on the proxy itself, the reset line, breaks the
coincidence while leaving the regime intact. That is what parasitism
means here. Every candidate depends on its environment. What matters is
whether its own macrostate carries the forecast or only stands in for
some off-partition variable that the observed process keeps aligned
with it. An organism kept alive by a stable environment is not parasitic
for that reason alone. It would be parasitic only if an admissible
perturbation acting on it showed that its apparent macro-dynamics had
tracked an external driver all along.

Both worries also clarify the countermodel's scope. It establishes a possibility, that the stated
observational criterion can certify a candidate whose autonomy is
borrowed, and it locates the condition under which this happens: some
admissible channel decouples the candidate from the structure sustaining
it. It shows nothing about how often that condition is met. Nothing here
suggests that parasitism is widespread in the special sciences, and the
argument does not need it to be. The conclusion is that the observational
criterion is insufficient in principle wherever such channels exist, and
whether they exist in a given domain is a question for that domain.

The proxy's clean macro-law holds only because,
in the observed process, \(Z_t = 0\) guarantees that the drivers agree.
The macro-variable does not itself track that agreement; the observed
process supplies it. Screening off
quantifies over the realizations the process happens to visit, and on
those realizations there is no heterogeneity to detect. Transition
closure quantifies over the realizations an admissible intervention can
reach, and there the heterogeneity is immediate. This is the same
distinction the exact benchmark already draws: the partition is weakly
lumpable, since the observed macro-process is genuinely Markov, but not
strongly lumpable, since strong lumpability demands equal onward
macro-transitions for every microstate in a class rather than only for
those the observed distribution happens to weight.

There is also a positive way to locate the autonomy the case contains.
The driver parity \(V_t = U^1_t \oplus U^2_t\) is itself a strongly
lumpable coarse-graining of the driver pair: since
\(V_{t+1} = V_t \oplus (N^1_{t+1} \oplus N^2_{t+1})\) whatever the reset
line does to the subsystems, every realization of \((U^1_t, U^2_t)\)
with a given parity has the same onward parity distribution, with
persistence \(r\). The autonomy in the case is therefore real; it is
just located at the driver partition rather than the candidate. \(Z\)
borrows \(V\)'s closure for as long as the dormant channel keeps the
two aligned. ``Parasitism'' is therefore a technical term here, not a metaphor. It names a candidate whose apparent transition structure is the closure of a different
partition, one the candidate neither contains nor preserves under
admissible perturbation.

The case is therefore a countermodel rather than a verbal trick or a
diagnostic failure: a candidate that satisfies unconditional relevance,
strict conditional irrelevance, and novelty, and that nonetheless fails
transition autonomy. Its role is to establish possibility, not
frequency. Nothing here claims that such candidates are common in the
special sciences, and the model illustrates a structural gap without
being evidence about any actual system. The gap itself is real: the
stated screening-off test cannot by itself exclude
borrowed autonomy, because the borrowing is invisible in the
distribution the test quantifies over.

The natural rejoinder, that screening off properly understood already
excludes the proxy, was met in \S\ref{why-closure-is-not-screening-off}:
the proxy is excluded on the modal reading, and closure is the test
that secures that reading.

\subsubsection{A Second Pressure Case: Hidden-Feedback Parasitism}
\label{hidden-feedback-case}

A second parasitism case arises when a macro-variable appears
transition-autonomous because a fast hidden process continually relaxes
it toward its apparent attractor. The source of the borrowed stability
is now internal and fast rather than upstream and exogenous, but the
skeleton of the previous case transfers directly, so the gap is not
tied to one configuration.

Retain the two subsystems \(A\) and \(B\) and the candidate
\(Z_t = A_t \oplus B_t\), and replace the external drivers with two fast
hidden variables \(H^1\) and \(H^2\) that relax toward the subsystems
they track. The micro-state is \(X_t = (A_t, B_t, H^1_t, H^2_t)\), and
the dynamics are that each subsystem's next value is its hidden
variable's current value, flipped with probability \(q\), after which
the hidden variables re-equilibrate to the subsystems. Because the
relaxation is fast relative to observation, the observed process never
leaves the equilibrated manifold on which \(H^1_t = A_t\) and
\(H^2_t = B_t\).

The verdicts are those of the previous case, and for the same reasons.
On the equilibrated manifold each subsystem performs an independent
symmetric walk with persistence \(1-q\), so \(Z\) is a symmetric chain
with persistence \(r = q^2 + (1-q)^2\), which differs from \(1-q\) and
cross-classifies the subsystem states exactly as before. Novelty is
therefore satisfied on Franklin and Robertson's condition rather than on
the weaker reading that a dependence is statable at the macro grain,
which is the reading their binary star already excludes. Within the
observed process, each macrostate is realized only by equilibrated
configurations, and those are homogeneous in their onward profile, so
conditional irrelevance holds strictly. An admissible perturbation that
acts faster than the relaxation sets a subsystem while its hidden
variable is still stale. That reaches a realization of the same
macrostate in which the hidden variables disagree, and from there
\(P(Z_{t+1} = 0) = 1 - r\) rather than \(r\). Closure fails.

The perturbation is admissible in the same three respects: it acts on
the candidate partition through a physically available channel, it
leaves the dynamics untouched, and the system re-equilibrates within a
step, so the regime is preserved and the deviation is a transient.

What differs between the two cases is the source of the borrowed
stability. The latent-driver proxy borrows from an upstream driver that
is exogenous and persistent; hidden-feedback parasitism borrows from an
internal relaxation that is fast and endogenous. The claim rests on
what they share. One structural fact generates the gap in both, that
the observed process visits only a homogeneous subset of a
heterogeneous macroclass, and the two sources show how little that fact
depends on where the off-partition structure sits. A defender of the
screening-off-plus-novelty package might concede the latent-driver proxy
as an edge phenomenon while keeping the criterion intact. Two sources
with the same signature are harder to relegate. Together they point to
a general category rather than a peculiar configuration: apparent
stability resting on structure the
partition itself does not preserve under admissible perturbation.

The intervention test here also differs from standard Woodwardian
causal testing. Standard interventionist analysis asks whether one
variable causes another within an already accepted partition. The test
here asks a prior question: whether a candidate partition itself
carries autonomous transition structure under admissible perturbation.

Without an intervention-indexed test, genuinely autonomous
macro-structure is indistinguishable from structure that merely
inherits apparent autonomy from upstream. This is the wedge that
motivates the rest of the paper, and it is why admissibility must be
declared before scoring rather than inferred from verdict patterns
afterward.

A similar point applies to causal-emergence work. Effective-information
analyses identify cases where macro-descriptions can be more causally
informative than micro-descriptions, often by improving effectiveness
through changes in determinism and degeneracy \citep{hoel2013};
\citet{dewhurst2021} argues that
such gains support only epistemic emergence absent further constraints.
That diagnosis is correct as far as it goes, but it identifies a
gap without filling it. A partition can increase effective information while remaining
fragile under intervention or horizon variation. The approach here
treats such gains as corroborating evidence within a closure-governed
criterion, not as a standalone ontology test.

Relatedly, Wallace's dynamic treatment of higher-level ontology is right
to push the debate away from mere collections of parts and toward
structure ``in the distributions and dynamics'' \citep{wallace2024}.
But a commuting or phase-space friendly representation is still not yet
a full admittance criterion. Closure asks when the higher-level
structure holds up under admissible perturbation rather than merely being
calibrated to observed trajectories. Ladyman's recent contribution to
the real-patterns volume moves in the same dynamic direction. He
characterizes real patterns in terms of ``projectible generalisations
that predict and explain'' \citep[p.~1]{ladyman2026}. The
present proposal is that his account still needs explicit admission
and downgrade discipline.

The paper therefore does not relabel computational mechanics or
causal-emergence analysis. It offers a philosophical criterion that
borrows their formal tools while adding intervention-indexed commitment conditions, an
explicit commitment protocol, and a sustained argument about why closure
is the right fix for pattern-realism's permissiveness problem.

\section{Non-Ideal Application: Leakiness, Diagnostics, and
  Verdicts}\label{approximate-closure-without-instrumentalist-drift}

\subsection{Why Approximation is
	Expected}\label{why-approximation-is-expected}

In complex systems, boundaries leak and couplings shift across regimes.
Exact closure is therefore unusual, and a credible realism criterion cannot
require perfect closure everywhere. Approximation is the expected non-ideal
form of the same criterion, provided constraints are fixed and diagnostics
are comparative. The study of approximate
aggregation has a long history, especially Simon's account of hierarchy
and near-decomposability, where strong within-subsystem couplings can
coexist with weaker cross-subsystem couplings and thereby separate
short-run local dynamics from longer-run aggregate behavior
\citep{simon1962,ando1961}. What this paper adds is philosophical. It
uses graded closure to support graded ontological verdicts and makes no
claim to improve estimation. Rosas et
al.'s distinction between informational, causal, and computational
closure is useful because it shows why the formal story is already
graded and plural at the diagnostic level even when used to motivate
one ontological criterion \citep{rosas2024}.

Ontology tests should be type-level before token-level. A pattern type
can be robustly closed for a regime even when a specific token fails
because of boundary violation, atypical perturbation, or timescale
mismatch. One anomalous token is therefore not decisive; what matters is
whether failures are exceptional or systematic for the type
under the declared constraints. Systematic failure, where within-class
micro-differences repeatedly matter for macro-transitions across tokens
under fixed constraints, triggers a type-level downgrade. Exceptional
failure is recorded as a token-level exception.

\subsection{Leakiness and the Distribution It Is Indexed
	To}\label{leakiness-as-canonical-target}

Leakiness measures how much within-class micro-detail still improves
macro-future prediction once the current macrostate is fixed. The
natural quantity is the conditional mutual information
\(I(X_t; Z_{t+1} \mid Z_t)\) \citep[eqs.~(2.60)--(2.61)]{cover2006}. In
plain terms, a non-zero value means that residual micro-detail still
helps predict the next macrostate even after the present one is fixed,
which is the signature of a leaky partition.

Stated that way the quantity is incomplete, and the gap affects the
criterion itself, not only estimation practice. A conditional
mutual information is defined relative to a joint distribution, and
nothing in the notation says which distribution over within-class
realizations is meant. Read against the observed joint, the measure
inherits exactly the limitation the paper's central argument turns on.
The latent-driver proxy of
\S\ref{what-the-formal-machinery-alone-does-not-provide} makes the point
sharply: there, screening off holds strictly, so
\(I(X_t; Z_{t+1} \mid Z_t)\) is exactly zero under the observed joint,
while the candidate fails transition autonomy. A measure that returns
perfect closure for a candidate the criterion must reject cannot be the
criterion's target quantity.

The repair is to index leakiness to the distribution over within-class
realizations. Doing so yields two quantities with different evidential
roles.

\begin{enumerate}
	\item
	      Observational leakiness: \(I(X_t; Z_{t+1} \mid Z_t)\) evaluated
	      under the observed joint. This is a screening diagnostic. It is
	      cheap, it requires no interventional access, and a
	      non-negligible value convicts immediately, since micro-detail
	      that demonstrably improves macro-prediction within the observed
	      process is transition-relevant heterogeneity the partition has
	      failed to absorb. What it cannot do is acquit. A value of zero
	      shows only that the realizations the process happens to visit
	      are homogeneous, which is consistent with borrowed autonomy.
	\item
	      Admissible leakiness: the same functional evaluated under the
	      distributions in which the within-class realization is set by an
	      admissible intervention. This is the criterion's quantity, and
	      its domain has to be stated, since the verdict turns on it. For
	      each macrostate \(z\), let \(R_I(z)\) be the microstates \(x\)
	      with \(g(x) = z\) that are reachable by interventions in the
	      admissible class \(I\) while preserving the regime, and let
	      \(\mathcal{M}_I(z)\) be the finite-support distributions over
	      \(R_I(z)\) induced by randomized mixtures of admissible
	      interventions. No single such distribution is privileged, so the
	      quantity is a supremum over them and over macrostates, taken at
	      the declared horizon \(L\):
	      \[
		      \Lambda^{(L)}_I(Z) \;=\; \sup_{z} \;\;
		      \sup_{\mu \in \mathcal{M}_I(z)}
		      I_\mu\!\left(X_t; Z_{t+1:t+L}\right),
	      \]
	      where conditioning on \(Z_t\) drops out because \(\mu\) is
	      already confined to one macroclass, and the target is the
	      \(L\)-step path rather than the next step alone, matching the
	      criterion of \S\ref{formal-statement}. The supremum
	      makes the quantity a test of the candidate. Under an average, a
	      partition could earn robust standing merely by being homogeneous on
	      some convenient weighting of its realizations.
\end{enumerate}

The kernels must be indexed to the interventions that produce them, not
to a surgery on the whole micro-state. Writing the pairwise condition
with \(do(X_t = x)\) would misdescribe the admissible class and would sit
badly with the restriction on Pearl-style surgery in
\S\ref{admissibility-and-hierarchical-evaluation}: an admissible
intervention here resets one subsystem and leaves the drivers untouched,
so it never sets \(X_t\) wholesale, and not every member of \(R_I(z)\)
need be reachable by setting the micro-state directly. Define instead the
post-intervention kernel
\[
	K_I\!\left(\cdot \mid x\right) \;=\;
	P\!\left(Z_{t+1:t+L} \in \cdot \;\middle|\;
	X_t = x,\ \text{admissible controls dormant thereafter}\right),
\]
for \(x \in R_I(z)\), and state closure as the requirement that
\(\sup_{z} \sup_{x, x' \in R_I(z)}
D\!\left(K_I(\cdot \mid x), K_I(\cdot \mid x')\right) = 0\)
for a divergence \(D\) separating distinct distributions.

The two formulations agree. If \(\mathcal{M}_I(z)\) contains every binary
mixture over \(R_I(z)\), then \(\Lambda^{(L)}_I(Z) = 0\) exactly when all
realizations in \(R_I(z)\) induce the same distribution over
\(Z_{t+1:t+L}\), for every \(z\). One direction is immediate: if the
kernels coincide on \(R_I(z)\), then \(Z_{t+1:t+L}\) is independent of
\(X_t\) under any \(\mu\) supported there, so the mutual information
vanishes. For the other, suppose \(K_I(\cdot \mid x) \neq
K_I(\cdot \mid x')\) for some \(x, x' \in R_I(z)\), and let \(\mu\) put
equal weight on the two. The conditional distribution of
\(Z_{t+1:t+L}\) then differs from its marginal on a set of positive
measure, so \(I_\mu(X_t; Z_{t+1:t+L}) > 0\) and the supremum is
positive. Either form therefore vanishes just in case within-class
micro-differences do not alter the macro-transition profile, which is
transition closure as stated in \S\ref{formal-statement}. In the
latent-driver proxy the observational quantity is exactly zero while both
admissible forms are strictly positive.

The two coincide whenever the observed process already visits the
reachable realizations of each macrostate with appreciable weight, which
is why observational diagnostics are so often adequate in practice and
why the literature has had no occasion to separate them. They come apart
exactly in the parasitism cases, where the observed distribution's
support is confined to a homogeneous subset of a heterogeneous class.

The same indexing shows where the exact benchmark fits. Strong lumpability, in Kemeny and
Snell's sense, requires that the coarse-grained process be Markov for
\emph{every} initial distribution, which is equivalent to requiring
equal onward macro-transitions for every microstate in a class
\citep{kemeny1960}. That is the limiting case in which admissible
leakiness vanishes for every distribution supported on a macroclass.
What observational leakiness can certify is weaker, and the contrast is
best put without leaning on a term whose formulations vary with the
initial distribution assumed: the observed process induces a Markov
macro-process for its actual initial or stationary distribution, while
the partition fails strong lumpability over the full microstate space.
The proxy is exactly such a case, and the two leakiness quantities
separate on it for that reason.

The asymmetry between the two is the same one recorded in
\S\ref{predictive-and-interventional-closure}: predictive adequacy can
persist where interventional stability fails, so predictive evidence is
important but not final. Observational leakiness remains the first thing
to compute and the cheapest way to eliminate candidates. It can screen
candidates out but cannot settle a verdict in their favor; a robust
verdict requires the admissible quantity.

Using information-theoretic measures to state an ontological criterion
might appear to assume a computationalist metaphysics. It does not. Both
quantities are functionals of transition kernels, and the conditional
independencies they measure hold or fail as facts about those kernels,
whatever anyone computes and whether or not any system in the world
exploits them (\S\ref{why-closure-carries-ontological-weight}). No claim
about the cost of tracking micro-detail, thermodynamic or otherwise, is
needed to make the measure objective.

Admissible leakiness is the target quantity, and observational leakiness
is its most accessible screen. Other diagnostics, such as within-class
transition divergence and predictive gain from added micro-features,
function as estimators or proxies when direct estimation is limited, and
they inherit the same indexing: what they are evidence for depends on
whether the within-class realizations they range over are those the
process visits or those an admissible intervention can reach.

Leakiness should therefore be read comparatively and through robustness
checks, not as a context-free absolute cutoff. Absolute tolerance is
domain-relative, but high sensitivity of leakiness rankings to modest
defensible modeling choices is itself a downgrade trigger.

Comparative ranking is necessary but not sufficient. A candidate
can beat nearby alternatives and still fail the criterion if absolute
leakiness remains high and instability persists under modest robustness
checks, so minimum viability is required as well as relative superiority.

The viability floor can be stated without a numeric threshold. A
partition fails minimum viability when, across the declared robustness
checks, within-class micro-features yield consistent, non-negligible
gains in macro-future prediction or intervention response relative to
the best macro-only model. If that pattern persists across
diagnostics and perturbation tests, the partition has not
achieved the transition autonomy the criterion requires, however it
ranks against competitors.

The term ``non-negligible'' can be given a comparative anchor without
fixing a universal numeric threshold. A within-class gain is
non-negligible when its magnitude is large enough that, were the same
gain observed between two candidate macro-partitions, it would be
recorded as a genuine comparative difference under the same diagnostic
conventions. The viability floor therefore measures the within-class
gain against a standard internal to the system, with no external
threshold. A partition fails if adding micro-features produces the kind
of effect size that would itself count as diagnostic evidence against a
rival candidate. Baseline model class and gain metric are fixed upstream
along with regime and admissibility under the canonical selection
discipline (\S\ref{procedural-safeguards-in-non-ideal-cases}). The
viability floor is applied first as a minimum gate; comparative
ranking then assigns verdict categories among candidates that pass it.

Approximate closure is therefore the non-ideal evidential route to the
same ontological target as exact closure, not a discretionary substitute
for it. Judgment is constrained by upstream fixing of
regime, admissibility, baseline model class, and gain metric, by
comparative stability across modest robustness checks, and by explicit
downgrade or indeterminacy when those rankings do not stabilize.

Verdict-category boundaries carry genuine vagueness because the
underlying property of transition autonomy is itself graded. Such
vagueness is no defect. Biological fitness is real and
explanatorily significant without admitting a sharp boundary between fit
and unfit. Ordering stability matters more than boundary placement. When partition
$A$ consistently shows lower leakiness and greater perturbation
stability than partition $B$ across diagnostics, that ordering should be
robust even if the exact placement of the boundary between ``qualified''
and ``robust'' carries some domain-conventional element. A partition
that dominates alternatives across diagnostics and remains stable under
modest changes to horizon, intervention distribution, and robustness
checks earns the stronger verdict regardless of where one draws the
categorical line.

\subsection{Procedural Safeguards in Non-Ideal
	Cases}\label{procedural-safeguards-in-non-ideal-cases}

Approximate closure claims require explicit safeguards. The canonical
selection discipline for the criterion, referenced throughout the paper,
is as follows.

\begin{enumerate}
	\item
	      Fix regime, horizon, admissibility class, diagnostics, baseline
	      model class, and gain metric before comparative scoring.
	\item
	      Compare candidate partitions under the same fixed setup.
	\item
	      Test robustness under modest changes in horizon and intervention
	      distribution.
	\item
	      Apply a minimum-viability floor: if all candidates remain high-leak
	      and fragile, return no robust standing verdict for that regime.
	\item
	      Report the closure profile and the evidential grade separately,
	      then the verdict they warrant
	      (\S\ref{reporting-and-evidential-discipline}).
	\item
	      Disclose any post hoc revision to the setup and restart the
	      evaluation rather than silently inheriting prior results.
\end{enumerate}

These safeguards also fit Wimsatt's robustness-analysis idea:
confidence should rise when partially independent ways of detecting the
same structure converge, and it should fall when apparent agreement is
driven by shared assumptions or disappears under small perturbation
\citep{wimsatt1981}. Multiple diagnostics do not replace criterion-level
argument with a vote count. They reduce the risk that a favorable verdict
is an artifact, or pseudorobust in Wimsatt's sense, of one estimator,
one horizon choice, or one admissibility specification.

The discipline keeps threshold choice procedural rather than
discretionary. It makes judgment auditable without eliminating it, and
it prevents favored partitions from being protected by moving criteria.
Difficulty of estimation is an epistemic limit on access and does not
show that the metaphysical criterion itself is defective. What is graded in non-ideal cases is not confidence alone but also the stability of closure itself across regime, horizon, admissibility, and robustness variation.

\subsection{Regime Dependence and
	Projectibility}\label{regime-dependence-and-projectibility}

Regime dependence is a claim about actual transition structure under
specified constraints, not about observer-relativity. A pattern can be closed in one regime and leaky in another
because the structure has changed, or because the available control
channels have shifted. Since the admissible intervention class is defined by the
regime's transition structure
(\S\ref{admissibility-and-hierarchical-evaluation}), a change in that
structure can alter both the admissible class and the closure verdict even when the
system's micro-dynamics remain the same. Regime and horizon index the
modal profile being claimed; specifying them does not license
arbitrary verdicts.

An ordinary physical analogy helps. Ice crystals are objective
structures only within certain temperature and pressure regimes. That
regime dependence shows that objective structure can be conditional on
equally objective background constraints without becoming epistemically
subjective.

Projectibility provides a further check. Defensible regime
specifications should support stable induction under modest
counterfactual variation. Narrowly engineered regimes that protect a
favored partition usually fail under small context shifts. When that
occurs, the candidate was never robustly closed in the relevant sense.

Here the paper is closest to Ladyman and Ross. Projectibility
functions as a realism-relevant stress test on whether a closure claim
survives beyond calibration conditions, and not as an optional
methodological virtue \citep{ladyman2007}. Closure without projectible
robustness is too weak a basis for robust realist claims.

Regime-indexing is also what makes disciplined robust verdicts possible
in the special sciences. Without explicit regime and intervention
indexing, most biological and social kinds would look merely leaky in an
undifferentiated way. With that indexing in place, the question becomes
whether the candidate is transition-autonomous within the bounded domain
that gives it its explanatory role.

Projectibility is still not sufficient by itself, because a pattern can
remain projectible while borrowing its stability from upstream structure
(see \S\ref{what-the-formal-machinery-alone-does-not-provide}). Closure
asks the further question of whether the candidate partition carries its
own transition structure under admissible perturbation.

Ladyman and Lorenzetti's effective ontic structural realism sharpens the
same point at the metaphysical level. If ontology is regime-dependent
and scale-relative, then what matters is whether a candidate is stable
enough in the relevant regime to merit commitment there, not whether it belongs to a single final inventory of being \citep{ladymanlorenzetti2024}.
Closure is offered here as the operational discipline for making that
judgment explicit.

Regime dependence itself should be split into two forms. Ontic regime
dependence occurs when system structure changes, such as phase
transitions or boundary-condition shifts. Epistemic regime dependence
occurs when measurement or control access changes while underlying
structure remains fixed. The first changes what is there to be tracked.
The second changes what can be warranted from available evidence.

\subsection{Failure Conditions and Downgrade
	Rules}\label{failure-conditions-and-downgrade-rules}

Commitment should be withdrawn or downgraded when any of four failure
patterns appears.

\begin{enumerate}
	\item
	      Persistent high leakiness across defensible diagnostics under fixed
	      constraints.
	\item
	      Strong disagreement among diagnostics that does not resolve under
	      modest robustness checks.
	\item
	      High sensitivity of verdicts to small, defensible changes in
	      admissibility or horizon.
	\item
	      Candidate sets where no partition clears the
	      minimum-viability floor.
\end{enumerate}

When these patterns appear, the right response is explicit downgrade to
qualified or indeterminate status, not a forced binary judgment. The
criterion does not require every case to produce a sharp ontology
verdict.

Conversely, when closure persists across distinct, independently
motivated admissibility classes within the same regime, that
cross-admissibility stability raises evidential confidence for robust
commitment.

As with the type-level tests above, a downgraded token does not
automatically refute type-level closure. What matters is whether instability is
exceptional at token level or systematic for the type under the stated
constraints.

\subsection{Stable Versus Merely Entailed
	Patterns}\label{stable-versus-merely-entailed-patterns}

A pattern can be entailed by microhistory and laws without being stable
as a macro-handle. Entailment is easy to come by; stability under
admissible perturbation is demanding. The distinction bears on
permissiveness debates. A contrived disjunctive construction can be true
of a realized trajectory while still failing closure, because it can
require continual within-class micro repair, fail under modest regime
shifts, and lose interventional reliability.

Phlogiston theory offers a historical illustration. It
compressed combustion phenomena under a single macro-variable, but the
partition required persistent bookkeeping to survive: when metals gained
weight upon burning, theorists introduced ``negative phlogiston''; when
different substances showed different weight changes, further ad hoc
parameters appeared. Each patch was a new within-class distinction
imported to preserve macro-prediction. Phlogiston was, in the
vocabulary of \S\ref{minimal-formal-illustration}, a high-maintenance
partition: reliable macro-prediction repeatedly forced refinement of
within-class distinctions. By contrast, oxygen theory required
substantially less such repair: the macro-variable (oxidation state)
screened off the relevant chemistry without persistent bookkeeping. The
criticism is structural and is not meant as retrospective ridicule. A candidate
that continually imports corrections to preserve macro-prediction
behaves like a non-autonomous partition, and the closure framework
detects this failure directly. The distinction between principled
theoretical development and ad hoc rescue is itself contested
philosophy of science, and the illustration is reconstructive rather
than a clean empirical test. Closure supplies one principled way to draw
the distinction, namely whether the
supplementation requires persistent within-class micro-bookkeeping or
introduces new transition-autonomous structure.

Gerrymandered constructions are therefore not necessarily false; they usually
fail to qualify for robust macro-level commitment. They may remain
descriptions, but not ontologically weighty descriptions in the relevant
regime. The same applies to computationally adequate models more
broadly. A higher-level model can be useful for prediction or control
without meeting closure standards. Computational adequacy means the
model serves some purpose; closure requires that transition-relevant
micro-differences are screened off in the specified regime. A
computationally adequate but high-leak representation remains a valuable
tool, but it does not automatically earn realist commitment. Closure
tests whether the relevant transition structure is in the world at the
target grain, rather than merely encoded in a successful representation.

\subsection{Practical Qualification Under Sparse Intervention
	Access}\label{practical-qualification-under-sparse-intervention-access}

In many domains, direct interventions are sparse, ethically constrained,
or expensive. Sparse intervention access changes how the criterion is
assessed; the criterion itself stays the same. The evidential order is
asymmetric. Predictive success can raise confidence, but only
interventional stability can underwrite robust commitment.

Where intervention data are limited, predictive closure functions as an
operational indicator while interventional closure remains the
ontological target. A model can fit historical trajectories and still
fail under novel perturbation. For that reason, claims should be
qualified by available intervention access and downgraded when
out-of-regime behavior is unstable. In-sample prediction is not
decisive. The standard remains robustness under admissible
perturbation, even when that robustness must be assessed indirectly.

Relevant evidence need not come only from humanly executable
interventions. In observational sciences it can also come from naturally
occurring disturbances, comparative cases, exogenous shocks, policy
changes, and other perturbations that bear on the same transition
structure. Such indirect intervention evidence often reveals whether a
macro-description remains stable when the system is pushed off its
ordinary trajectory.

The same considerations favor triangulating diagnostics.
Predictive closure can look strong in-sample while interventional
closure fails under admissible perturbation. Partial observability can
hide within-class heterogeneity that later appears as instability.
Nonstationarity can make a previously low-leak partition unstable
outside its calibration window. Triangulation does not remove these
risks, but it makes them visible.

Payment systems offer a useful illustration, since direct intervention
on them is rare and mostly not ours to perform. A well-functioning
payment system can look observationally stable for long stretches. What
matters for stronger commitment is how that stability behaves under
settlement delays, enforcement shocks, trust disruptions, or policy
changes that function as admissible perturbations. Where those
perturbations leave transactional organization intact, robust verdicts
become plausible. Where the same structure survives only
observationally, or collapses once such pressures appear, the right
verdict is qualified rather than robust.

\subsection{Verdict Rule and Reporting
	Discipline}\label{reporting-and-evidential-discipline}

Two questions bear on a verdict, and they can come apart. Whether a
candidate is transition-autonomous is a question about the candidate.
Whether we have shown that it is autonomous is a question about our
evidence. A verdict is what the two together warrant.

\begin{enumerate}
	\item
	      \emph{Closure profile}, a property of the candidate under the
	      declared regime, horizon, and admissible intervention class.
	      \emph{Closed} when admissible leakiness is at or near zero;
	      \emph{partially closed} when it is low within the declared
	      constraints but rises under modest variation of them;
	      \emph{high-leak} when within-class micro-differences materially
	      alter the macro-transition profile. Nothing here is about what
	      anyone knows.
	\item
	      \emph{Evidential grade}, a property of our access to that
	      profile. \emph{Strong} when interventional evidence spans the
	      admissible class and diagnostics converge; \emph{limited} when
	      the evidence is largely predictive because intervention access is
	      sparse, costly, or ethically closed
	      (\S\ref{practical-qualification-under-sparse-intervention-access});
	      \emph{unstable} when diagnostics disagree persistently or
	      rank-orders do not settle.
	\item
	      \emph{Commitment verdict}, the attitude the first two warrant.
	      \emph{Robust} for a closed profile established on strong
	      evidence. \emph{Qualified} for a partially closed profile on
	      strong evidence, or a closed profile on limited evidence, since
	      fragility in the candidate and thinness in the evidence both
	      justify commitment short of robust, for different reasons.
	      \emph{Indeterminate} when the evidential grade is unstable,
	      whatever the profile may be. The label describes our epistemic
	      position. What is indeterminate is the verdict, not the
	      transition structure it is about.
\end{enumerate}

The separation also brings out a case that a single verdict scale would
hide. A high-leak profile on strong evidence yields a determinate
negative verdict. The candidate has been shown to lack transition
autonomy, and robust standing is denied rather than deferred. Treating
that as indeterminacy would report a finding as a failure to find.
Conversely, an indeterminate verdict says nothing about the candidate. A
candidate can be perfectly closed and still receive one, and better
evidence may later show that it was.

The rule sets a floor. Stricter domain conventions can narrow the robust
category further, but nothing weaker than this triangulation should
underwrite robust verdicts.

Where the canonical selection discipline
(\S\ref{procedural-safeguards-in-non-ideal-cases}) fixes \emph{what is
	held constant} before scoring, reporting fixes \emph{what is disclosed}
once scoring is done. To keep conclusions comparable across cases,
closure assessments should state five items explicitly: target
partition and rationale; regime and horizon specification; admissible
intervention class and justification; diagnostics used and their
agreement or disagreement; and final verdict category. Stating these
prevents silent shifts in target, timescale, or admissibility from being
mistaken for genuine ontological progress.

Reports should also give the profile and the grade alongside the
verdict, since the verdict alone does not say which of them drove it. A
qualified verdict on a fragile candidate and a qualified verdict on a
well-behaved candidate we have barely probed are different situations
with different remedies. The first calls for narrowing the regime, the
second for better evidence. For the same reason, the anti-instrumentalist
claim of \S\ref{why-closure-carries-ontological-weight} attaches to the
profile rather than to the verdict. The profile is a fact about the
candidate; the verdict is an attitude, answerable to evidence as well as
to the world.

Diagnostics do not replace criterion-level argument. They provide
comparative evidence about whether a candidate satisfies the criterion
under the constraints fixed for the assessment. No single proxy is
treated as infallible. Three checks are usually sufficient in practice:
a leakiness proxy based on predictive gain from added within-class
micro-features; within-class transition-profile divergence; and
intervention-response invariance under admissible perturbation.
Agreement across these checks increases warrant; persistent disagreement
is evidence for revision or downgrade rather than for forced commitment.
When interventional evidence is sparse or unavailable, the remaining
diagnostics can still support qualified status, but robust standing
requires that the interventional leg not be entirely missing.

\section{Objections and Replies}\label{objections-and-replies}

Two objections are taken here. A third, that closure merely renames
screening off properly understood, was answered in
\S\ref{why-closure-is-not-screening-off}, where the concession is made
directly. Three standing commitments apply to both replies and are not
restated in each. First, the account does not deny microphysical
completeness; it claims non-redundancy at the explanatory and
interventional level for declared macro-targets, so the causal-exclusion
pressure is met by target-indexing rather than by multiplying causes.
Second, all verdicts are indexed to the declared constraints, so
objections presupposing context-free ontological claims address a view
the paper does not hold. Third, the strongest conclusion is conditional
on structural realist and interventionist commitments; readers who
withhold those commitments can still accept the narrower
anti-gerrymandering result.

\subsection{``Instrumentalism and
	Admissibility''}\label{instrumentalism-and-admissibility}

Objection: target selection is interest-relative, so the view still
collapses into usefulness.

Reply: target selection can be interest-shaped without making closure
verdicts preference-shaped. Once regime, horizon, and admissibility are
fixed, whether macro-transitions are autonomous is a system fact under
explicit constraints.

The stronger version invokes Dennett's Martian
\citep[p.~42]{dennett1991}: if an ideal micro-predictor can forecast
everything, macro realism looks optional. But this infers ontological
redundancy from representational power. Even a micro-omniscient
predictor coexists with objectively better macro-bottlenecks for target
dynamics, and ignoring those bottlenecks increases representational
burden without erasing macro transition structure. A more powerful
observer may have less use for the macro-description, but whether a
partition is transition-autonomous under fixed constraints does not vary
with who computes it.

The admissibility version of the same objection holds that
admissibility criteria already encode what counts as the relevant
level, so the circularity worry returns one level up. In reply,
admissibility selection is constrained by the system's actual control
ecology, meaning the physically realizable, regime-preserving
perturbations the system already faces under thermal fluctuation,
diffusion, mechanical contact, and similar naturally occurring sources.
The analyst's task is to identify that ecology, not to construct it.
Hierarchical order and disclosure rules make the identification
auditable: admissibility is fixed from physical constraints before
partition scoring, and post hoc admissibility revision triggers restart.
Where the ecology is genuinely contested, the criterion correctly
returns indeterminate verdicts rather than forcing a favored outcome.
The related worry, that admissibility merely ``shifts'' the hard
problem, assumes that the choice of intervention class is free, whereas
the framework treats it as a physically anchored feature of the regime
being evaluated. A problem is merely shifted when it is replaced by an
equally hard one. Here it is replaced by a problem that the regime
itself at least partly settles, since which control channels exist is
not something the analyst decides. The choice is not thereby automatic,
and the account does not claim that it is.

Residual disagreement can remain, and the account does not promise
universal convergence from one pass. What it offers is disciplined
comparison and explicit downgrade when verdicts are unstable across
nearby defensible admissibility specifications, so that hard cases are
flagged as unresolved rather than settled by stipulation.

Instrumentalist pressure takes two forms: ontology should track
predictive utility only, or ontology is unnecessary once predictive
utility is available. Both fail. Utility without closure is too
permissive, and closure without ontology understates what stable
intervention-guiding structure provides. What closure adds is a
discipline governing commitment, without metaphysical extravagance. If a
candidate repeatedly supports stable counterfactual control under fixed
constraints, withholding ontological standing is a position that owes an
account of what further condition realism was supposed to require.
Rather than forcing maximal realism, the paper asks for explicit criteria
for when anti-realism remains credible.

The objection that leakiness is ``only'' predictive error has real force
against the observational quantity, which is why that quantity screens
rather than settles (\S\ref{leakiness-as-canonical-target}). Directed at
admissible leakiness, the objection presupposes a gap between tracking
causal structure and describing what is real. But under the paper's
structural realist commitments, there is no further ontological layer
for that gap to open onto. Stable counterfactual and predictive
organization is not evidence for some deeper condition left unspecified
elsewhere; it is what the relevant reality consists in at that grain.
When a partition leaks under admissible intervention, then, the
candidate lacks autonomous transition structure, and there is no real
object hiding underneath that we failed to see. Admissible leakiness
therefore measures ontic deficiency in the candidate partition, not
merely our ignorance.

The cheap-coding objection fails for the same reason. A recoding can
improve local fit and still fail closure under admissible perturbation,
so the criterion cannot be satisfied by representational choice alone.
That resistance to cheap recoding is what permissiveness critiques
demand.

\subsection{``Closure Is Too Strong and Eliminates Most Special
	Sciences''}\label{closure-is-too-strong-and-eliminates-most-special-sciences}

Objection: if closure requires that within-class micro-differences not
matter for macro-transitions, then most special-science kinds will fail.
Biological kinds are notoriously leaky \citep{boyd1991}, and many
psychological states and economic categories are no sharper. The criterion eliminates the
ontology it was supposed to discipline.

Reply: the objection conflates robust standing with any positive verdict
at all. Three verdict categories are available. Robust standing requires
stable closure under fixed constraints. Qualified standing fits cases
where closure is partial, regime-limited, or supported by convergent but
incomplete evidence. Indeterminate status fits cases where evidence is
too unstable to warrant commitment.

Most special-science kinds fall into the qualified range, and that is
the correct result. Biological species exhibit substantial transition
autonomy within ecological and evolutionary regimes while leaking at
boundary cases (hybrid zones, ring species, horizontal gene transfer).
Psychological kinds often support predictive and interventional
regularity within clinical or experimental regimes while failing under
broader perturbation. Rather than eliminating these kinds, the criterion
gives them the status their closure evidence supports: qualified
standing with explicit regime limitations, short of full robust status.

Nor must exact transition autonomy hold everywhere for a
special-science kind to count as robustly real. Because the admissible
intervention class helps define the regime, a macro-variable such as a
biological trait or an economic variable can achieve robust standing
within a narrow but physically bounded ecological, physiological, or
institutional home even if it leaks outside that home. It need only be
transition-autonomous under the intervention class that defines the
regime in which it operates. Downgrade to qualified status becomes
appropriate when those regime boundaries are stretched, or when
intervention evidence within the regime remains too sparse to warrant
the stronger verdict.

Nearby real-pattern work in biology provides a useful foil. One can
characterize biological objects by their explanatory power and pragmatic
utility, and by the availability of multiple valid carvings, without yet
having shown that lower-level variation is transition-irrelevant in the
stronger sense used here \citep{jiang2025}. Such utility matters, but it
does not by itself settle robust standing.

A sharper form of the objection concedes the graded verdicts and
targets their distribution: the criterion is decisive where it is
easy, in declared-channel settings like the Markov benchmark, and
returns qualified or indeterminate verdicts almost everywhere the
debate matters, so its bite is inversely proportional to its
applicability. Three points blunt this. First, qualified verdicts are comparative results, not silence. They rank candidate partitions, locate where
leakiness lives, and say what evidence would upgrade the verdict, which
is more than a permissive criterion says anywhere. Second, rival
criteria do not perform better where closure verdicts are hard; they
issue robust verdicts there, and the wedge cases show that some of those
verdicts certify borrowed autonomy. A criterion that withholds robust
standing where the evidence cannot support it is not weaker than one
that grants it there. Third, decisiveness tracks the maturity of a
domain's control ecology. As intervention practice accumulates, the same
criterion delivers stronger verdicts, so its reach is not confined to
clean cases by construction.

The objection therefore proves too much. A criterion that granted robust
standing to every special-science kind regardless of leakiness would
be the permissive view closure is designed to replace. Graded verdicts
let the assessment match the available evidence instead of forcing a
binary choice between full realism and eliminativism.

\section{Conclusion: A Disciplined Realism
  Claim}\label{conclusion-a-disciplined-realism-claim}

The three claims distinguished at the outset rest on different things.
The anti-gerrymandering constraint rests on the wedge cases of
\S\ref{what-the-formal-machinery-alone-does-not-provide} and requires no
particular realism package. The criterion for robust commitment rests on
the structural-realist and interventionist commitments of
\S\ref{intervention-structure-and-the-causation-worry} and
\S\ref{why-closure-carries-ontological-weight}. The
structured-space reading rests on neither, and is offered as a direction
rather than a result: a reader who grants the first two can withhold it
without loss to the admittance-and-downgrade protocol.

Rainforest realism is compelling, but its admittance criteria remain
incomplete. Screening off plus novelty remains necessary, but
observational success can be parasitic. The latent-driver proxy and
hidden-feedback cases
(\S\ref{what-the-formal-machinery-alone-does-not-provide},
\S\ref{hidden-feedback-case}) show that a candidate can satisfy
observational screening off while borrowing its apparent autonomy from
off-partition organization the partition itself does not preserve
under admissible perturbation. Closure under admissible intervention
is what distinguishes autonomous from borrowed macro-structure. Strong
lumpability provides the exact benchmark in Markov settings. Leakiness
and convergent diagnostics extend the same criterion to non-ideal
settings as a graded verdict protocol with explicit downgrade
discipline.

This further step sharpens the non-redundancy demand already implicit in
the Primacy of Physics Constraint
(\S\ref{closure-and-the-primacy-of-physics-constraint}): a candidate is
non-redundant when lower-level variation within its macroclasses stops
mattering to the relevant macro-transition profile under the declared
constraints. In the recent literature, the criterion supplies what
Franklin and Robertson's screening-off test, Rosas et al.'s formal
diagnostics, Ladyman and Lorenzetti's effective realism, and recent
dynamic formulations of pattern realism do not yet by themselves
supply: an explicit closure-based admittance and downgrade discipline
for robust macro-level commitment
\citep{franklinrobertson2024,rosas2024,ladymanlorenzetti2024,wallace2024,ladyman2026}.

The proposal supports selective commitment: robust where closure is
stable, qualified in borderline cases, explicitly downgraded where
verdicts depend on fragile admissibility or horizon choices. The overall
position is intentionally middle-range. It is stronger than
compression-only pattern realism because it adds exclusion and downgrade
structure, and weaker than a total-level metaphysics because it stays
restricted to induced closure in spatiotemporal systems under declared constraints.

On the programmatic reading of rainforest realism as a structured space
of computationally closed coarse-grainings, the verdict categories make the
structure reportable: robust where placement is stable, qualified where
ontological support is real but fragile, and indeterminate where no
stable placement has yet been secured. The formal richness of that
space in any given domain depends on what the closure discipline turns
up there, and the core argument does not depend on any particular such
finding.

The view separates itself from eliminativist pressure by neither
treating higher-level descriptions as fictions by default nor granting
them standing by convenience alone. The test is whether they track
stable transition structure under admissible intervention. If they do,
realist commitment follows on the paper's stated terms. If they do not,
realism should be downgraded.

This paper does not offer a complete metaphysics of levels: the scope stays restricted
to induced closure in spatiotemporal systems under declared
constraints. It does not settle general debates about emergence or
reduction; closure is one necessary filter within a cumulative
admittance protocol, not a standalone theory of higher-level reality.
It supplies the dynamical component of an account of objecthood and
leaves room for complementary accounts of how closed structure is
realized, classified, and individuated (\S\ref{scope-and-non-claims}).
It does not supply a universal implementation recipe: applications
beyond the illustrations depend on domain-specific estimator choices,
robustness checks, and control-channel inventories that remain
downstream of the philosophical argument and do not alter its content.
And the strongest conclusion depends on the structural realist and
interventionist commitments stated in
\S\ref{intervention-structure-and-the-causation-worry} and
\S\ref{why-closure-carries-ontological-weight}; a reader who withholds
those commitments can still accept the narrower anti-gerrymandering
result. What the paper contributes is a philosophical criterion together
with a discipline for applying it.

Future work should test the staged admittance protocol across domains
by comparing candidate partitions under fixed constraints, rather than
by multiplying new concepts. Worries about circularity, instrumentalist
drift, and scope ambiguity have been addressed directly, so the replies
can be judged on the merits of the stated argument. One further
question, left for future work, is whether this closure-based account
of realist commitment bears on substrate-independence or on more
abstract forms of structural standing.

\backmatter

\label{references}
\bibliography{real-patterns-need-closure-SYNTHESE}

\end{document}
